\frfilename{mt251.tex}
\versiondate{10.11.06}
\copyrightdate{2000}

\def\chaptername{Ukuran produk}
\def\sectionname{Produk berhingga}

\newsection{251}

Konstruksi pertama yang perlu disusun adalah produk dari sepasang ruang
ukur. Ternyata, bahkan untuk mencari metode kanonis yang dapat diterapkan
secara universal, sudah terdapat kesulitan teknis yang cukup besar. Karena
itu, saya perlu menguraikan dua konstruksi yang berkaitan tetapi berbeda,
yaitu ukuran produk `primitif' dan `c.l.d.' (251C, 251F). Setelah
mencantumkan sifat-sifat mendasar ukuran produk c.l.d\ (251I-251J), saya
menguraikan identifikasi produk ukuran Lebesgue dengan dirinya sendiri
(251N) dan pembahasan yang cukup menyeluruh tentang subruang (251O-251S).

\vleader{48pt}{251A}{Definisi} Misalkan $(X,\Sigma,\mu)$ dan $(Y,\Tau,\nu)$
dua
ruang ukur. Untuk $A\subseteq X\times Y$, tetapkan

\Centerline{$\theta A=\inf\{\sum_{n=0}^{\infty}\mu E_n\cdot\nu F_n:
  E_n\in\Sigma,\,F_n\in\Tau\,\Forall \,n\in\Bbb N,
  \,A\subseteq\bigcup_{n\in\Bbb N}E_n\times F_n\}$.}

\cmmnt{\medskip

\noindent{\bf Catatan} Dalam hasil kali $\mu E_n\cdot\nu F_n$,
$0\cdot\infty$ harus diartikan sebagai $0$, seperti dalam \S135.
}

\leader{251B}{Lema} Dalam konteks 251A, $\theta$ merupakan ukuran luar
pada $X\times Y$.

\proof{{\bf (a)} Dengan menetapkan $E_n=F_n=\emptyset$ untuk setiap
$n\in\Bbb N$, kita melihat bahwa $\theta\emptyset = 0$.

\medskip

{\bf (b)} Jika $A\subseteq B\subseteq X\times Y$, maka setiap kali
$B\subseteq\bigcup_{n\in\Bbb N}E_n\times F_n$, kita juga mempunyai
$A\subseteq\bigcup_{n\in\Bbb N}E_n\times F_n$; jadi $\theta A\le\theta B$.

\medskip

{\bf (c)} Misalkan $\sequencen{A_n}$ suatu barisan subhimpunan $X\times Y$
dengan gabungan $A$.
Untuk setiap $\epsilon>0$, bagi tiap $n\in\Bbb N$ kita dapat memilih barisan
$\sequence{m}{E_{nm}}$ dalam $\Sigma$ dan $\sequence{m}{F_{nm}}$ dalam $\Tau$
sedemikian sehingga $A_n\subseteq\bigcup_{m\in\Bbb N}E_{nm}\times F_{nm}$ dan
$\sum_{m=0}^{\infty}\mu E_{nm}\cdot\nu F_{nm}
\le\theta A_n+2^{-n}\epsilon$.
Karena $\Bbb N\times\Bbb N$ terhitung, terdapat bijeksi
$k\mapsto(n_k,m_k):\Bbb N\to\Bbb N\times\Bbb N$, dan sekarang

\Centerline{$A\subseteq\bigcup_{n,m\in\Bbb N}E_{nm}\times F_{nm}
=\bigcup_{k\in\Bbb N}E_{n_km_k}\times F_{n_km_k}$,}

\noindent sehingga

$$\eqalign{\theta A&\le\sum_{k=0}^{\infty}
  \mu E_{n_km_k}\cdot\nu F_{n_km_k}
=\sum_{n=0}^{\infty}\sum_{m=0}^{\infty}\mu E_{nm}\cdot\nu F_{nm}\cr
&\le\sum_{n=0}^{\infty}\theta A_n+2^{-n}\epsilon
=2\epsilon+\sum_{n=0}^{\infty}\theta A_n.\cr}$$

\noindent Karena $\epsilon$ sebarang, $\theta
A\le\sum_{n=0}^{\infty}\theta A_n$.

Karena $\sequencen{A_n}$ sebarang, $\theta$ merupakan ukuran luar.
}

\leader{251C}{Definisi} Misalkan $(X,\Sigma,\mu)$ dan $(Y,\Tau,\nu)$
ruang ukur. Yang saya maksud dengan {\bf ukuran produk primitif} pada
$X\times Y$ adalah ukuran $\lambda_0$ yang diperoleh melalui metode
\Caratheodory\cmmnt{ (113C)} dari ukuran luar $\theta$ yang didefinisikan
dalam 251A.

\cmmnt{\medskip

\noindent{\bf Catatan}
Perlu saya kemukakan bahwa tidak ada kesepakatan umum mengenai ukuran mana
yang seharusnya disebut sebagai `ukuran produk' pada $X\times Y$. Bahkan, dalam
251F di bawah ini saya akan memperkenalkan suatu ukuran alternatif, dan dalam
catatan untuk bagian ini saya akan menyebutkan ukuran yang ketiga.
}

\leader{251D}{Definisi}\cmmnt{ Akan memudahkan jika suatu konstruksi alami
bagi aljabar-$\sigma$ diberi nama.} Jika $X$ dan $Y$ adalah himpunan dengan
aljabar-$\sigma$ $\Sigma\subseteq\Cal PX$ dan $\Tau\subseteq\Cal PY$, saya akan
menuliskan $\Sigma\tensorhat\Tau$ untuk aljabar-$\sigma$ subhimpunan
$X\times Y$ yang dibangkitkan oleh
$\{E\times F:E\in\Sigma,\,F\in\Tau\}$.


\leader{251E}{Proposisi} Misalkan $(X,\Sigma,\mu)$ dan $(Y,\Tau,\nu)$
ruang ukur; misalkan $\lambda_0$ ukuran produk primitif pada $X\times Y$,
dan $\Lambda$ domainnya.
Maka $\Sigma\tensorhat\Tau\subseteq\Lambda$ dan
$\lambda_0(E\times F)=\mu E\cdot\nu F$ untuk semua $E\in\Sigma$
dan $F\in\Tau$.

\proof{{\bf (a)} Misalkan $E\in\Sigma$ dan $A\subseteq X\times Y$.
Untuk setiap $\epsilon>0$, terdapat barisan $\sequencen{E_n}$ dalam $\Sigma$
dan $\sequencen{F_n}$ dalam $\Tau$ sedemikian sehingga
$A\subseteq\bigcup_{n\in\Bbb N}E_n\times F_n$ dan
$\sum_{n=0}^{\infty}\mu E_n\cdot\nu F_n\le\theta A+\epsilon$. Sekarang

\Centerline{$A\cap(E\times Y)
\subseteq\bigcup_{n\in\Bbb N}(E_n\cap E)\times F_n$,
\quad$A\setminus(E\times Y)
\subseteq\bigcup_{n\in\Bbb N}(E_n\setminus E)\times F_n$,}

\noindent sehingga

$$\eqalign{\theta(A\cap(E\times Y))+\theta(A\setminus(E\times Y))
&\le\sum_{n=0}^{\infty}\mu(E_n\cap E)\cdot\nu F_n
  +\sum_{n=0}^{\infty}\mu(E_n\setminus E)\cdot\nu F_n\cr
&=\sum_{n=0}^{\infty}\mu E_n\cdot\nu F_n
\le\theta A+\epsilon.\cr}$$

\noindent Karena $\epsilon$ sebarang,
$\theta(A\cap(E\times Y))+\theta(A\setminus(E\times Y))\le\theta A$.
Dan ini cukup untuk memastikan bahwa $E\times Y\in\Lambda$ (lihat 113D).

\medskip

{\bf (b)} Dengan cara serupa, $X\times F\in\Lambda$ untuk setiap $F\in\Tau$,
sehingga $E\times F=(E\times Y)\cap(X\times F)\in\Lambda$ untuk setiap
$E\in\Sigma$, $F\in\Tau$.

Karena $\Lambda$ merupakan aljabar-$\sigma$, ia harus memuat aljabar-$\sigma$
terkecil yang memuat semua produk $E\times F$, yaitu
$\Lambda\supseteq\Sigma\tensorhat\Tau$.

\medskip


{\bf (c)} Ambil $E\in\Sigma$, $F\in\Tau$. Kita mengetahui bahwa $E\times
F\in\Lambda$; dengan menetapkan $E_0=E$, $F_0=F$, $E_n=F_n=\emptyset$ untuk
$n\ge 1$ dalam definisi $\theta$, kita memperoleh

\Centerline{$\lambda_0(E\times F)
=\theta(E\times F)\le\mu E\cdot\nu F$.}

Kita telah sampai pada gagasan utama konstruksi ini. Sesungguhnya,
$\theta(E\times F)=\mu E\cdot\nu F$. \Prf\ Misalkan
$E\times F\subseteq\bigcup_{n\in\Bbb N}E_n\times F_n$, dengan
$E_n\in\Sigma$ dan $F_n\in\Tau$ untuk setiap $n$. Tetapkan
$u=\sum_{n=0}^{\infty}\mu E_n\cdot\nu F_n$. Jika $u=\infty$ atau $\mu E=0$
atau $\nu F=0$, tentu saja $\mu E\cdot\nu F\le u$. Jika tidak, tetapkan

\Centerline{$I=\{n:n\in\Bbb N,\,\mu E_n=0\}$,
\quad$J=\{n:n\in\Bbb N,\,\nu F_n=0\}$,
\quad$K=\Bbb N\setminus(I\cup J)$,}

\Centerline{$E'=E\setminus\bigcup_{n\in I}E_n$,
\quad$F'=F\setminus\bigcup_{n\in J}F_n$.}

\noindent Maka $\mu E'=\mu E$ dan $\nu F'=\nu F$;
$E'\times F'\subseteq\bigcup_{n\in K}E_n\times F_n$; dan untuk $n\in K$,
$\mu E_n<\infty$ dan $\nu F_n<\infty$, karena
$\mu E_n\cdot\nu F_n\le u<\infty$ dan
baik $\mu E_n$ maupun $\nu F_n$ tidak nol. Tetapkan

\Centerline{$f_n=\nu F_n\chi E_n:X\to\Bbb R$}

\noindent jika $n\in K$, dan $f_n=\tbf{0}:X\to\Bbb R$ jika $n\in I\cup J$.
Maka $f_n$ suatu fungsi sederhana dan
$\int f_n=\nu F_n\mu E_n$ untuk $n\in K$, serta $0$ jika tidak, sehingga

\Centerline{$\sum_{n=0}^{\infty}$
$\textfont3=\twelveex\int f_n(x)\mu(dx)$
$=\sum_{n=0}^{\infty}\mu E_n\cdot\nu F_n
\le u$.}

\noindent Menurut teorema B.Levi (123A), yang diterapkan pada
$\sequencen{\sum_{k=0}^nf_k}$,
$g=\sum_{n=0}^{\infty}f_n$ terintegralkan dan $\int g\,d\mu\le u$.
Tuliskan $E''$ untuk $\{x:x\in E',\,g(x)<\infty\}$, sehingga
$\mu E''=\mu E'=\mu E$. Sekarang ambil sebarang $x\in E''$ dan tetapkan
$K_x=\{n:n\in K,\,x\in E_n\}$. Karena
$E'\times F'\subseteq\bigcup_{n\in K}E_n\times F_n$,
$F'\subseteq\bigcup_{n\in K_x}F_n$ dan

\Centerline{$\nu F=\nu F'\le\sum_{n\in K_x}\nu F_n
=\sum_{n=0}^{\infty}f_n(x)=g(x)$.}

\noindent Jadi $g(x)\ge\nu F$ untuk setiap $x\in E''$. Kita
mengasumsikan bahwa $0<\mu E=\mu E''$ dan $0<\nu F$, sehingga haruslah
$\nu F<\infty$, $\mu E''<\infty$. Sekarang $g\ge\nu F\chi E''$, sehingga

\Centerline{$\mu E\cdot\nu F=\mu E''\cdot\nu F
=\int\nu F\chi E''\le\int g\le u$
$=\sum_{n=0}^{\infty}\mu E_n\cdot\nu F_n$.}

\noindent Karena $\sequencen{E_n}$, $\sequencen{F_n}$ sebarang,
$\theta(E\times F)\ge\mu E\cdot\nu F$ dan
$\theta(E\times F)=\mu E\cdot\nu F$.\ \Qed

Jadi

\Centerline{$\lambda_0(E\times F)=\theta(E\times F)=\mu E\cdot\nu F$}

\noindent untuk semua $E\in\Sigma$, $F\in\Tau$.
}%end of proof of 251E

\leaveitout{untuk (c), M.Burke mengusulkan:

$\chi E(x)\chi F(y)\le\sum_{n=0}^{\infty}\chi E_n(x)\chi F_n(y)$

menurut 135F-135G, dengan mengintegralkan terhadap $y$, \query

$\chi E(x)\nu F\le\sum_{n=0}^{\infty}\chi E_n(x)\nu F_n$

dan sekarang

$\mu E\cdot\nu F\le\sum_{n=0}^{\infty}\mu E_n\cdot\nu F_n$.}


\leader{251F}{Definisi} Misalkan $(X,\Sigma,\mu)$ dan $(Y,\Tau,\nu)$
ruang ukur, dan $\lambda_0$ ukuran produk primitif
\cmmnt{ yang didefinisikan dalam 251C}.
Yang saya maksud dengan {\bf ukuran produk c.l.d.\ } pada $X\times Y$ adalah
fungsi $\lambda:\dom\lambda_0\to[0,\infty]$ yang didefinisikan dengan

\Centerline{$\lambda W=\sup\{\lambda_0(W\cap(E\times F)):E\in\Sigma,\,
F\in\Tau,\,\mu E<\infty,\,\nu F<\infty\}$}

\noindent untuk $W\in\dom\lambda_0$.

\leader{251G}{Catatan}\cmmnt{ Sebaiknya segera saya tunjukkan
bahwa} $\lambda$ merupakan suatu
ukuran. \prooflet{\Prf\ Tentu saja domainnya $\Lambda=\dom\lambda_0$
merupakan aljabar-$\sigma$,
dan $\lambda\emptyset=\lambda_0\emptyset=0$. Jika $\sequencen{W_n}$ suatu
barisan saling lepas dalam $\Lambda$, maka untuk setiap
$E\in\Sigma$, $F\in\Tau$ yang berukuran hingga,

\Centerline{$\lambda_0(\bigcup_{n\in\Bbb N}W_n\cap(E\times F))
=\sum_{n=0}^{\infty}\lambda_0(W_n\cap(E\times F))
\le\sum_{n=0}^{\infty}\lambda W_n$,}

\noindent sehingga
$\lambda(\bigcup_{n\in\Bbb N}W_n)\le\sum_{n=0}^{\infty}\lambda W_n$.
Di pihak lain, jika
$a<\sum_{n=0}^{\infty}\lambda W_n$, kita dapat menemukan $m\in\Bbb N$ dan
$a_0,\ldots,a_m$ sedemikian sehingga $a\le\sum_{n=0}^ma_n$ dan
$a_n<\lambda W_n$ untuk setiap $n\le m$; sekarang terdapat
$E_0,\ldots,E_m\in\Sigma$ dan $F_0,\ldots,F_m\in\Tau$, semuanya berukuran
hingga, sedemikian sehingga
$a_n\le\lambda_0(W_n\cap(E_n\times F_n))$ untuk setiap $n$. Dengan menetapkan
$E=\bigcup_{n\le m}E_n$ dan $F=\bigcup_{n\le m}F_n$, kita mempunyai
$\mu E<\infty$ dan $\nu F<\infty$, sehingga
$$\eqalign{\lambda(\bigcup_{n\in\Bbb N}W_n)
&\ge\lambda_0(\bigcup_{n\in\Bbb N}W_n\cap(E\times F))
=\sum_{n=0}^{\infty}\lambda_0(W_n\cap(E\times F))\cr
&\ge\sum_{n=0}^m\lambda_0(W_n\cap(E_n\times F_n))
\ge\sum_{n=0}^ma_n
\ge a.\cr}$$

\noindent Karena $a$ sebarang, $\lambda(\bigcup_{n\in\Bbb
N}W_n)\ge\sum_{n=0}^{\infty}\lambda W_n$ dan $\lambda(\bigcup_{n\in\Bbb
N}W_n)=\sum_{n=0}^{\infty}\lambda W_n$. Karena $\sequencen{W_n}$
sebarang, $\lambda$ merupakan suatu ukuran.\ \Qed}

\leader{251H}{}\cmmnt{ Kita memerlukan suatu sifat sederhana dari ukuran
$\lambda_0$.

\medskip

\noindent}{\bf Lema} Misalkan $(X,\Sigma,\mu)$ dan $(Y,\Tau,\nu)$ dua
ruang ukur; misalkan $\lambda_0$ ukuran produk primitif pada
$X\times Y$, dan $\Lambda$ domainnya. Jika $H\subseteq X\times Y$ dan
$H\cap(E\times F)\in\Lambda$ setiap kali
$\mu E<\infty$ dan $\nu F<\infty$, maka $H\in\Lambda$.

\proof{ Misalkan $\theta$ ukuran luar yang diuraikan dalam 251A. Misalkan
$A\subseteq X\times Y$ dan $\theta A<\infty$. Ambil
$\epsilon>0$. Misalkan $\sequencen{E_n}$, $\sequencen{F_n}$ berturut-turut
barisan dalam $\Sigma$, $\Tau$ sedemikian sehingga
$A\subseteq\bigcup_{n\in\Bbb N}E_n\times F_n$ dan
$\sum_{n=0}^{\infty}\mu E_n\cdot\nu F_n\le\theta A+\epsilon$.
Sekarang, untuk setiap $n$, hasil kali ukuran $\mu E_n$, $\nu F_n$
berhingga, sehingga salah satunya nol atau
keduanya berhingga. Jika $\mu E_n=0$ atau $\nu F_n=0$, tentu saja

\Centerline{$\mu E_n\cdot\nu F_n=0
=\theta((E_n\times F_n)\cap H)+\theta((E_n\times F_n)\setminus H)$.}

\noindent Jika $\mu E_n<\infty$ dan $\nu F_n<\infty$, maka

$$\eqalignno{\mu E_n\cdot\nu F_n
&=\lambda_0(E_n\times F_n)\cr
&=\lambda_0((E_n\times F_n)\cap H)+\lambda_0((E_n\times F_n)\setminus
H)\cr
&=\theta((E_n\times F_n)\cap H)+\theta((E_n\times F_n)\setminus H).\cr
}$$

\noindent Oleh karena itu, karena $\theta$ suatu ukuran luar,

$$\eqalign{\theta(A\cap H)+\theta(A\setminus H)
&\le\sum_{n=0}^{\infty}\theta((E_n\times F_n)\cap H)
  +\sum_{n=0}^{\infty}\theta((E_n\times F_n)\setminus H)\cr
&=\sum_{n=0}^{\infty}\mu E_n\cdot\nu F_n
\le\theta A+\epsilon.\cr}$$

\noindent Karena $\epsilon$ sebarang, $\theta(A\cap H)+\theta
(A\setminus H)\le\theta A$. Karena $A$ sebarang, $H\in\Lambda$.
}

\leader{251I}{}\cmmnt{ Sekarang untuk sifat-sifat mendasar dari
ukuran produk c.l.d.

\medskip

\noindent}{\bf Teorema} Misalkan $(X,\Sigma,\mu)$ dan $(Y,\Tau,\nu)$
ruang ukur; misalkan $\lambda$ ukuran produk c.l.d. pada
$X\times Y$, dan $\Lambda$ domainnya. Maka

(a) $\Sigma\tensorhat\Tau\subseteq\Lambda$ dan
$\lambda(E\times F)=\mu E\cdot\nu F$ setiap kali $E\in\Sigma$, $F\in\Tau$
dan $\mu E\cdot\nu F<\infty$;

(b) untuk setiap $W\in\Lambda$ terdapat $V\in\Sigma\tensorhat\Tau$
sedemikian sehingga $V\subseteq W$ dan $\lambda V=\lambda W$;

(c) $(X\times Y,\Lambda,\lambda)$ lengkap dan ditentukan secara lokal,
dan bahkan merupakan versi c.l.d. dari
$(X\times Y,\Lambda,\lambda_0)$\cmmnt{ sebagaimana diuraikan dalam
213D-213E}; khususnya, $\lambda W=\lambda_0W$ setiap kali
$\lambda_0W<\infty$;

(d) jika $W\in\Lambda$ dan $\lambda W>0$, maka terdapat $E\in\Sigma$,
$F\in\Tau$ sedemikian sehingga $\mu E<\infty$, $\nu F<\infty$ dan
$\lambda(W\cap(E\times F))>0$;

(e) jika $W\in\Lambda$ dan $\lambda W<\infty$, maka untuk setiap
$\epsilon>0$ terdapat $E_0,\ldots,E_n\in\Sigma$,
$F_0,\ldots,F_n\in\Tau$, semuanya berukuran hingga, sedemikian sehingga
$\lambda(W\symmdiff\bigcup_{i\le n}(E_i\times F_i))\le\epsilon$.

\proof{ Ambil $\theta$ sebagai ukuran luar dari 251A dan
$\lambda_0$ sebagai ukuran produk primitif dari 251C. Tetapkan
$\Sigma^f=\{E:E\in\Sigma,\,\mu E<\infty\}$ dan
$\Tau^f=\{F:F\in\Tau,\,\nu F<\infty\}$.

\medskip

{\bf (a)} Menurut 251E, $\Sigma\tensorhat\Tau\subseteq\Lambda$. Jika
$E\in\Sigma$ dan $F\in\Tau$ serta $\mu E\cdot\nu F<\infty$, maka salah satu
kemungkinan adalah $\mu E\cdot\nu F=0$ dan
$\lambda(E\times F)=\lambda_0(E\times F)=0$, atau
$\mu E$ dan $\nu F$ keduanya berhingga dan kembali
$\lambda(E\times F)=\lambda_0(E\times F)=\mu E\cdot\nu F$.

\medskip

{\bf (b)(i)} Ambil sebarang $a<\lambda W$. Maka terdapat
$E\in\Sigma^f$, $F\in\Tau^f$ sedemikian sehingga
$\lambda_0(W\cap(E\times F))>a$ (251F); sekarang

$$\eqalign{\theta((E\times F)\setminus W)
&=\lambda_0((E\times F)\setminus W)\cr
&=\lambda_0(E\times F)-\lambda_0(W\cap(E\times F))
<\lambda_0(E\times F)-a.\cr}$$

\noindent Misalkan $\sequencen{E_n}$, $\sequencen{F_n}$ berturut-turut
barisan dalam $\Sigma$, $\Tau$ sedemikian sehingga
$(E\times F)\setminus W\subseteq\bigcup_{n\in\Bbb N}E_n\times F_n$ dan
$\sum_{n=0}^{\infty}\mu E_n\cdot\nu F_n
\le\lambda_0(E\times F)-a$. Pertimbangkan

\Centerline{$V=(E\times
F)\setminus\bigcup_{n\in\Bbb N}E_n\times F_n\in\Sigma\tensorhat\Tau$;}

\noindent maka $V\subseteq W$, dan

$$\eqalignno{\lambda V=\lambda_0V
&=\lambda_0(E\times F)-\lambda_0((E\times F)\setminus V)\cr
&\ge\lambda_0(E\times F)-\lambda_0(\bigcup_{n\in\Bbb N}E_n\times F_n)\cr
\noalign{\noindent (karena $(E\times F)\setminus
V\subseteq\bigcup_{n\in\Bbb N}E_n\times F_n$)}
&\ge\lambda_0(E\times F)-\sum_{n=0}^{\infty}\mu E_n\cdot\nu F_n
\ge a\cr}$$

\noindent (berdasarkan pemilihan $E_n$, $F_n$).

\medskip

\quad{\bf (ii)} Jadi, untuk setiap $a<\lambda W$ terdapat
$V\in\Sigma\tensorhat\Tau$ sedemikian sehingga $V\subseteq W$ dan
$\lambda V\ge a$. Sekarang pilih barisan $\sequencen{a_n}$ yang meningkat
tegas menuju $\lambda W$, dan bagi setiap $a_n$ pilih $V_n$ yang
bersesuaian; maka $V=\bigcup_{n\in\Bbb N}V_n$ termasuk dalam aljabar-$\sigma$
$\Sigma\tensorhat\Tau$, termuat dalam $W$, dan mempunyai ukuran sekurang-
kurangnya $\sup_{n\in\Bbb N}\lambda V_n$ dan paling banyak $\lambda W$;
jadi $\lambda V=\lambda W$, sebagaimana diperlukan.

\medskip

{\bf (c)(i)} Jika $H\subseteq X\times Y$ terabaikan terhadap $\lambda$,
terdapat $W\in\Lambda$ sedemikian sehingga $H\subseteq W$ dan
$\lambda W=0$. Jika $E\in\Sigma$, $F\in\Tau$ berukuran hingga,
$\lambda_0(W\cap(E\times F))=0$; tetapi karena $\lambda_0$ diperoleh dari
ukuran luar $\theta$
dengan metode \Caratheodory, ia lengkap (212A), sehingga
$H\cap(E\times F)\in\Lambda$ dan $\lambda_0(H\cap(E\times F))=0$.
Karena $E$ dan $F$
sebarang, menurut 251H kita mempunyai $H\in\Lambda$. Karena $H$ sebarang,
$\lambda$ lengkap.

\medskip

\quad{\bf (ii)} Jika $W\in\Lambda$ dan $\lambda W=\infty$, maka harus
terdapat $E\in\Sigma$, $F\in\Tau$ sedemikian sehingga $\mu E<\infty$,
$\nu F<\infty$ dan $\lambda_0(W\cap(E\times F))>0$; sekarang

\Centerline{$0<\lambda(W\cap(E\times F))\le\mu E\cdot\nu F<\infty$.}

\noindent Jadi $\lambda$ semihingga.

\medskip

\quad{\bf (iii)} Jika $H\subseteq X\times Y$ dan $H\cap W\in\Lambda$
setiap kali $\lambda W<\infty$, maka, khususnya, $H\cap(E\times
F)\in\Lambda$ setiap kali $\mu E<\infty$ dan $\nu F<\infty$; kembali
menurut 251H, $H\in\Lambda$. Jadi $\lambda$ ditentukan secara lokal.

\medskip

\quad{\bf (iv)} Jika $W\in\Lambda$ dan $\lambda_0W<\infty$, terdapat
barisan $\sequencen{E_n}$ dalam $\Sigma$, $\sequencen{F_n}$ dalam $\Tau$
sedemikian
sehingga $W\subseteq\bigcup_{n\in\Bbb N}(E_n\times F_n)$ dan
$\sum_{n=0}^{\infty}\mu E_n\cdot\nu F_n<\infty$. Tetapkan

\Centerline{$I=\{n:\mu E_n=\infty\}$,
\quad$J=\{n:\nu F_n=\infty\}$,
\quad$K=\Bbb N\setminus (I\cup J)$;}

\noindent maka $\nu(\bigcup_{n\in I}F_n)=\mu(\bigcup_{n\in J}E_n)=0$,
sehingga $\lambda_0(W\setminus W')=0$, dengan

\Centerline{$W'=W\cap\bigcup_{n\in K}(E_n\times F_n)
\supseteq W\setminus((\bigcup_{n\in J}E_n\times Y)
  \cup(X\times\bigcup_{n\in I}F_n))$.}

\noindent Sekarang tetapkan $E'_n=\bigcup_{i\in K,i\le n}E_i$,
$F'_n=\bigcup_{i\in K,i\le n}F_i$ untuk setiap $n$. Kita mempunyai
$W'=\bigcup_{n\in\Bbb N}W'\cap(E'_n\times F'_n)$, sehingga

\Centerline{$\lambda
W\le\lambda_0W=\lambda_0W'=\lim_{n\to\infty}\lambda_0(W'\cap(E'_n\times
F'_n))
\le \lambda W'
\le\lambda W$,}

\noindent dan $\lambda W=\lambda_0W$.

\medskip

\quad{\bf (v)} Mengikuti terminologi 213D, kita tuliskan

\Centerline{$\tilde\Lambda=\{W:W\subseteq X\times Y,\,W\cap V\in\Lambda$
setiap kali $V\in\Lambda$ dan $\lambda_0V<\infty\}$,}

\Centerline{$\tilde\lambda W
=\sup\{\lambda_0(W\cap V):V\in\Lambda$, $\lambda_0V<\infty\}$.}

\noindent Karena $\lambda_0(E\times F)<\infty$ setiap kali
$\mu E<\infty$ dan $\nu F<\infty$, $\tilde\Lambda\subseteq\Lambda$ dan
$\tilde\Lambda=\Lambda$.


Sekarang, untuk setiap $W\in\Lambda$ kita mempunyai

$$\eqalignno{\tilde\lambda W
&=\sup\{\lambda_0(W\cap V):V\in\Lambda,\,\lambda_0V<\infty\}\cr
&\ge\sup\{\lambda_0(W\cap (E\times F)):E\in\Sigma^f,\,F\in\Tau^f\}\cr
&=\lambda W\cr
&\ge\sup\{\lambda(W\cap V):V\in\Lambda,\,\lambda_0V<\infty\}\cr
&=\sup\{\lambda_0(W\cap V):V\in\Lambda,\,\lambda_0V<\infty\},\cr}$$

\noindent dengan menggunakan (iv) tepat di atas, sehingga
$\lambda=\tilde\lambda$ merupakan versi c.l.d. dari $\lambda_0$.

\medskip

{\bf (d)} Jika $W\in\Lambda$ dan $\lambda W>0$, terdapat
$E\in\Sigma^f$ dan $F\in\Tau^f$ sedemikian sehingga
$\lambda(W\cap(E\times F))=\lambda_0(W\cap(E\times F))>0$.

\medskip

{\bf (e)} Terdapat $E\in\Sigma^f$, $F\in\Tau^f$ sedemikian sehingga
$\lambda_0(W\cap(E\times F))\ge\lambda W-\bover13\epsilon$; tetapkan
$V_1=W\cap (E\times F)$; maka

\Centerline{$\lambda(W\setminus V_1)=\lambda W-\lambda V_1
=\lambda W-\lambda_0V_1\le\Bover13\epsilon$.}

\noindent Terdapat
barisan $\sequencen{E'_n}$ dalam $\Sigma$, $\sequencen{F'_n}$ dalam
$\Tau$ sedemikian sehingga $V_1\subseteq\bigcup_{n\in\Bbb
N}E'_n\times F'_n$ dan $\sum_{n=0}^{\infty}\mu E'_n\cdot\nu
F'_n\le\lambda_0V_1+\bover13\epsilon$. Dengan mengganti $E'_n$, $F'_n$
dengan $E'_n\cap E$, $F'_n\cap F$ jika perlu, kita dapat menganggap bahwa
$E'_n\in\Sigma^f$ dan $F'_n\in\Tau^f$ untuk setiap $n$. Tetapkan
$V_2=\bigcup_{n\in\Bbb N}E'_n\times F'_n$; maka

\Centerline{$\lambda(V_2\setminus V_1)
\le\lambda_0(V_2\setminus V_1)
\le\sum_{n=0}^{\infty}\mu E'_n\cdot\nu F'_n-\lambda_0V_1
\le\Bover13\epsilon$.}

\noindent Misalkan $m\in\Bbb N$ sedemikian sehingga
$\sum_{n=m+1}^{\infty}\mu E'_n\cdot\nu F'_n\le\bover13\epsilon$, dan
tetapkan

\Centerline{$V=\bigcup_{n=0}^mE'_n\times F'_n$.}

\noindent Maka

\Centerline{$\lambda(V_2\setminus V)
\le\sum_{n=m+1}^{\infty}\mu E'_n\cdot\nu F'_n\le\Bover13\epsilon$.}

Dengan menggabungkan semua ini, kita memperoleh
$W\symmdiff V\subseteq(W\setminus V_1)\cup(V_2\setminus
V_1)\cup(V_2\setminus V)$, sehingga

\Centerline{$\lambda(W\symmdiff V)
\le\lambda(W\setminus V_1)+\lambda(V_2\setminus
V_1)+\lambda(V_2\setminus
V)\le\Bover13\epsilon+\Bover13\epsilon+\Bover13\epsilon=\epsilon$.}

\noindent Dan $V$ mempunyai bentuk yang disyaratkan.
}

\leader{251J}{Proposisi} Jika $(X,\Sigma,\mu)$ dan $(Y,\Tau,\nu)$
ruang ukur semihingga dan $\lambda$ ukuran produk c.l.d. pada
$X\times Y$, maka $\lambda(E\times F)=\mu E\cdot\nu F$ untuk semua
$E\in\Sigma$, $F\in\Tau$.

\proof{ Dengan menetapkan $\Sigma^f=\{E:E\in\Sigma,\,\mu E<\infty\}$,
$\Tau^f=\{F:F\in\Tau,\,\nu F<\infty\}$, kita mempunyai

$$\eqalign{\lambda(E\times F)
&=\sup\{\lambda_0((E\cap E_0)\times(F\cap F_0)):
  E_0\in\Sigma^f,\,F_0\in\Tau^f\}\cr
&=\sup\{\mu(E\cap E_0)\cdot\nu(F\cap F_0)):
  E_0\in\Sigma^f,\,F_0\in\Tau^f\}\cr
&=\sup\{\mu(E\cap E_0):E_0\in\Sigma^f\}\cdot
  \sup\{\nu(F\cap F_0):F_0\in\Tau^f\}
=\mu E\cdot\nu F\cr}$$

\noindent (dengan menggunakan 213A).
}%end of proof of 251J
\leader{251K}{$\sigma$-hingga \dvrocolon{ruang}}\cmmnt{ Tentu saja
sebagian besar ruang ukur yang padanya hasil-hasil ini akan kita terapkan
bersifat $\sigma$-hingga, dan dalam kasus ini terdapat beberapa penyederhanaan yang berguna.

\medskip

\noindent}{\bf Proposisi} Misalkan $(X,\Sigma,\mu)$ dan $(Y,\Tau,\nu)$
ruang ukur $\sigma$-hingga. Maka ukuran produk c.l.d.\ pada
$X\times Y$ sama dengan ukuran produk primitif, dan merupakan
kelengkapan dari pembatasannya pada $\Sigma\tensorhat\Tau$; selain itu, ukuran
produk yang sama ini bersifat $\sigma$-hingga.

\proof{ Tuliskan $\lambda_0$, $\lambda$ untuk ukuran produk primitif dan c.l.d.\
seperti biasa, serta $\Lambda$ untuk domain keduanya. Misalkan
$\sequencen{E_n}$, $\sequencen{F_n}$ barisan tak menurun dari himpunan-himpunan
berukuran hingga yang masing-masing menutupi $X$, $Y$ (lihat 211D).

\medskip

{\bf (a)} Untuk setiap $n\in\Bbb N$, $\lambda(E_n\times F_n)=\mu E_n\cdot\nu
F_n$ berhingga, menurut 251Ia. Karena $X\times Y=\bigcup_{n\in\Bbb
N}E_n\times F_n$, $\lambda$ bersifat $\sigma$-hingga.

\medskip

{\bf (b)} Untuk sebarang $W\in\Lambda$,

\Centerline{$\lambda_0W=\lim_{n\to\infty}\lambda_0(W\cap(E_n\times
F_n))=\lim_{n\to\infty}\lambda(W\cap(E_n\times F_n))=\lambda W$.}

\noindent Jadi $\lambda=\lambda_0$.

\medskip

{\bf (c)} Tuliskan $\lambda_{\Cal B}$ untuk pembatasan
$\lambda=\lambda_0$ pada $\Sigma\tensorhat\Tau$, dan
$\hat\lambda_{\Cal B}$ untuk kelengkapannya.

\medskip

\quad{\bf (i)} Misalkan $W\in\dom\hat\lambda_{\Cal B}$. Maka terdapat
$W'$, $W''\in\Sigma\tensorhat\Tau$ sedemikian sehingga
$W'\subseteq W\subseteq W''$ dan $\lambda_{\Cal B}(W''\setminus W')=0$
(212C). Dalam hal ini,
$\lambda(W''\setminus W')=0$; karena $\lambda$ lengkap, $W\in\Lambda$
dan

\Centerline{$\lambda W=\lambda W'=\lambda_{\Cal B}W'
=\hat\lambda_{\Cal B}W$.}

\noindent Jadi $\lambda$ memperluas $\hat\lambda_{\Cal B}$.

\medskip

\quad{\bf (ii)} Jika $W\in\Lambda$, maka terdapat
$V\in\Sigma\tensorhat\Tau$ sedemikian sehingga $V\subseteq W$ dan
$\lambda(W\setminus V)=0$. \Prf\ Untuk setiap $n\in\Bbb N$ terdapat
$V_n\in\Sigma\tensorhat\Tau$ sedemikian sehingga
$V_n\subseteq W\cap(E_n\times F_n)$ dan
$\lambda V_n=\lambda(W\cap(E_n\times F_n))$ (251Ib). Tetapi karena
$\lambda(E_n\times F_n)=\mu E_n\cdot\nu F_n$ berhingga, ini berarti
$\lambda(W\cap(E_n\times F_n)\setminus V_n)=0$. Jadi jika kita menetapkan
$V=\bigcup_{n\in\Bbb N}V_n$, kita memperoleh $V\in\Sigma\tensorhat\Tau$,
$V\subseteq W$ dan

\Centerline{$W\setminus V
=\bigcup_{n\in\Bbb N}W\cap(E_n\times F_n)\setminus V
\subseteq\bigcup_{n\in\Bbb N}W\cap(E_n\times F_n)\setminus V_n$}

\noindent terabaikan terhadap $\lambda$.\ \Qed

Demikian pula, terdapat $V'\in\Sigma\tensorhat\Tau$ sedemikian sehingga
$V'\subseteq(X\times Y)\setminus W$ dan
$\lambda(((X\times Y)\setminus W)\setminus V')=0$. Dengan menetapkan $V''=(X\times Y)\setminus V'$,
$V''\in\Sigma\tensorhat\Tau$, $W\subseteq V''$ dan
$\lambda(V''\setminus W)=0$. Jadi

\Centerline{$\lambda_{\Cal B}(V''\setminus V)
=\lambda(V''\setminus V)
=\lambda(V''\setminus W)+\lambda(W\setminus V)=0$,}

\noindent dan $W$ terukur terhadap $\hat\lambda_{\Cal B}$, dengan
$\hat\lambda_{\Cal B}W=\lambda_{\Cal B}V=\lambda W$. Karena $W$
sembarang, $\hat\lambda_{\Cal B}=\lambda$.
}%end of proof of 251K

\leader{*251L}{}\cmmnt{ Hasil berikut sesuai ditempatkan di sini;
saya memberinya tanda bintang karena hasil ini jarang digunakan (terdapat versi yang lebih penting
dari gagasan yang sama dalam 254G) dan buktinya
dapat dilewati sampai Anda memerlukannya.

\medskip

\noindent}{\bf Proposisi}
Misalkan $(X_1,\Sigma_1,\mu_1)$, $(X_2,\Sigma_2,\mu_2)$,
$(Y_1,\Tau_1,\nu_1)$ dan $(Y_2,\Tau_2,\nu_2)$ ruang ukur
$\sigma$-hingga; misalkan
$\lambda_1$, $\lambda_2$ ukuran produk masing-masing pada $X_1\times Y_1$,
$X_2\times Y_2$. Andaikan $f:X_1\to X_2$ dan
$g:Y_1\to Y_2$ adalah fungsi-fungsi \imp, dan
$h(x,y)=(f(x),g(y))$ untuk $x\in X_1$, $y\in Y_1$. Maka $h$ bersifat \imp.
%251K

\proof{ Tuliskan $\Lambda_1$, $\Lambda_2$ untuk domain dari $\lambda_1$,
$\lambda_2$ masing-masing.

\medskip

{\bf (a)} Misalkan $E\in\Sigma_2$ dan $F\in\Tau_2$ berukuran hingga.
Maka $\lambda_1h^{-1}[W\cap(E\times F)]$ terdefinisi dan sama dengan
$\lambda_2(W\cap(E\times F))$ untuk setiap $W\in\Lambda_2$. \Prf\

$$\eqalign{\lambda_1h^{-1}[E\times F]
&=\lambda_1(f^{-1}[E]\times g^{-1}[F])
=\mu_1f^{-1}[E]\cdot\nu_1g^{-1}[F]\cr
&=\mu_2E\cdot\nu_2F
=\lambda_2(E\times F)\cr}$$

\noindent menurut 251E/251J.\ \Qed

\medskip

{\bf (b)} Ambil $E_0\in\Sigma_2$ dan $F_0\in\Tau_2$ berukuran hingga.
Misalkan $\tilde\lambda_1$, $\tilde\lambda_2$ ukuran subruang masing-masing pada
$f^{-1}[E_0]\times g^{-1}[F_0]$ dan $E_0\times F_0$.
Tetapkan $\tilde h=h\restr f^{-1}[E_0]\times g^{-1}[F_0]$, dan tuliskan
$\iota$ untuk pemetaan identitas dari $E_0\times F_0$ ke $X_2\times Y_2$; misalkan
$\lambda=\tilde\lambda_1\tilde h^{-1}$ dan
$\lambda'=\tilde\lambda_2\iota^{-1}$ merupakan
ukuran citra pada $X_2\times Y_2$. Maka (a) memberi tahu kita bahwa

$$\eqalign{\lambda(E\times F)
&=\lambda_1(h^{-1}[(E\cap E_0)\times(F\cap F_0)])\cr
&=\lambda_2((E\cap E_0)\times(F\cap F_0))
=\lambda'(E\times F)\cr}$$

\noindent setiap kali $E\in\Sigma_2$ dan $F\in\Tau_2$. Menurut Teorema Kelas
Monoton (136C), $\lambda$ dan $\lambda'$ berimpit pada
$\Sigma_2\tensorhat\Tau_2$, yaitu,
$\lambda_1(h^{-1}[W\cap(E_0\times F_0)])=\lambda_2(W\cap(E_0\times F_0))$
untuk setiap $W\in\Sigma_2\tensorhat\Tau_2$.

Jika $W$ sebarang anggota $\Lambda_2$,
terdapat $W'$, $W''\in\Sigma_2\tensorhat\Tau_2$ sedemikian
sehingga $W'\subseteq W\subseteq W''$ dan $\lambda_2(W''\setminus W')=0$
(251K). Sekarang haruslah

\Centerline{$h^{-1}[W'\cap(E_0\times F_0)]
\subseteq h^{-1}[W\cap(E_0\times F_0)]
\subseteq h^{-1}[W''\cap(E_0\times F_0)]$,}

\Centerline{$\lambda_1(h^{-1}[W''\cap(E_0\times F_0)]\setminus
h^{-1}[W'\cap(E_0\times F_0)])
=\lambda_2((W''\setminus W')\cap(E_0\times F_0))=0$;}

\noindent karena $\lambda_1$ lengkap,
$\lambda_1h^{-1}[W\cap(E_0\times F_0)]$ terdefinisi dan sama dengan

\Centerline{$\lambda_1h^{-1}[W'\cap(E_0\times F_0)]
=\lambda_2(W'\cap(E_0\times F_0))
=\lambda_2(W\cap(E_0\times F_0))$.}

\medskip

{\bf (c)} Sekarang misalkan $\sequencen{E_n}$, $\sequencen{F_n}$ merupakan
barisan tak menurun dari himpunan-himpunan berukuran hingga dengan gabungan masing-masing $X_2$, $Y_2$.
Jika $W\in\Lambda_2$,

\Centerline{$\lambda_1h^{-1}[W]
=\sup_{n\in\Bbb N}\lambda_1h^{-1}[W\cap(E_n\times F_n)]
=\sup_{n\in\Bbb N}\lambda_2(W\cap(E_n\times F_n))
=\lambda_2W$.}

\noindent Jadi $h$ bersifat \imp, sebagaimana dinyatakan.
}%end of proof of 251L

\leader{251M}{}\cmmnt{ Sudah waktunya saya memberikan beberapa contoh. Tentu saja contoh
utamanya
adalah ukuran Lebesgue. Dalam hal ini kita memperoleh satu-satunya hasil yang wajar.
Saya berhenti sejenak untuk menguraikan contoh utama dari produk
$\Sigma\tensorhat\Tau$ yang diperkenalkan dalam 251D.

\medskip

\noindent}{\bf Proposisi} Misalkan $r$, $s\ge 1$ bilangan bulat. Maka kita
mempunyai bijeksi alami $\phi:\BbbR^r\times\BbbR^s\to\BbbR^{r+s}$,
yang didefinisikan dengan menetapkan

\Centerline{$\phi((\xi_1,\ldots,\xi_r),(\eta_1,\ldots,\eta_s))
=(\xi_1,\ldots,\xi_r,\eta_1,\ldots,\eta_s)$}

\noindent untuk $\xi_1,\ldots,\xi_r,\eta_1,\ldots,\eta_s\in\Bbb R$. Jika
kita tuliskan $\Cal B_r$, $\Cal B_s$ dan $\Cal B_{r+s}$ untuk aljabar-sigma Borel
dari $\BbbR^r$, $\BbbR^s$ dan $\BbbR^{r+s}$
masing-masing, maka $\phi$ mengidentifikasi
$\Cal B_{r+s}$ dengan $\Cal B_r\tensorhat\Cal B_s$.

\proof{{\bf (a)} Tuliskan $\Cal B$ untuk aljabar-sigma
$\{\phi^{-1}[W]:W\in\Cal B_{r+s}\}$ yang disalin ke $\BbbR^r\times\BbbR^s$
melalui bijeksi $\phi$; yang hendak kita buktikan adalah
$\Cal B=\Cal B_r\tensorhat\Cal B_s$. Kita mempunyai pemetaan
$\pi_1:\BbbR^{r+s}\to\BbbR^r$, $\pi_2:\BbbR^{r+s}\to\BbbR^s$ yang didefinisikan dengan
menetapkan $\pi_1(\phi(x,y))=x$, $\pi_2(\phi(x,y))=y$. Setiap koordinat
dari $\pi_1$ kontinu, sehingga terukur Borel (121Db), jadi
$\pi_1^{-1}[E]\in\Cal B_{r+s}$ untuk setiap $E\in\Cal B_r$, menurut 121K.
Demikian pula,
$\pi_2^{-1}[F]\in\Cal B_{r+s}$ untuk setiap $F\in\Cal B_s$. Jadi
$\phi[E\times F]=\pi_1^{-1}[E]\cap\pi_2^{-1}[F]$ termasuk dalam
$\Cal B_{r+s}$, yakni, $E\times F\in\Cal B$, setiap kali $E\in\Cal B_r$
dan $F\in\Cal B_s$. Karena $\Cal B$ merupakan
aljabar-sigma, $\Cal B_r\tensorhat\Cal B_s\subseteq\Cal B$.

\medskip

{\bf (b)} Sekarang periksa himpunan-himpunan berbentuk

\Centerline{$\{(x,y):\xi_i\le\alpha\}
=\{x:\xi_i\le\alpha\}\times\BbbR^s$,}

\Centerline{$\{(x,y):\eta_j\le\alpha\}
=\BbbR^r\times\{y:\eta_j\le\alpha\}$}

\noindent untuk $\alpha\in\Bbb R$, $i\le r$ dan $j\le s$, dengan mengambil
$x=(\xi_1,\ldots,\xi_r)$ dan $y=(\eta_1,\ldots,\eta_s)$. Semua himpunan ini
termasuk dalam $\Cal B_r\tensorhat\Cal B_s$. Tetapi aljabar-sigma yang dibangkitkannya
tepat sama dengan $\Cal B$, menurut 121J. Jadi
$\Cal B\subseteq\Cal B_r\tensorhat\Cal B_s$ dan
$\Cal B=\Cal B_r\tensorhat\Cal B_s$.
}%end of proof of 251M

\leader{251N}{Teorema} Misalkan $r$, $s\ge 1$ bilangan bulat. Maka
bijeksi $\phi:\BbbR^r\times\BbbR^s\to\BbbR^{r+s}$ yang diuraikan dalam 251M
mengidentifikasi ukuran Lebesgue pada $\BbbR^{r+s}$ dengan
produk c.l.d.\ $\lambda$
dari ukuran Lebesgue pada $\BbbR^r$ dan ukuran Lebesgue pada $\BbbR^s$.

\proof{ Tuliskan $\mu_r$,
$\mu_s$, $\mu_{r+s}$ untuk ketiga versi ukuran Lebesgue,
$\mu^*_r$, $\mu^*_s$ dan $\mu^*_{r+s}$ untuk ukuran luar yang bersesuaian,
dan
$\theta$ untuk ukuran luar pada $\BbbR^r\times\BbbR^s$ yang diperoleh dari
$\mu_r$ dan $\mu_s$ melalui rumus dalam 251A.

\medskip

{\bf (a)} Jika $I\subseteq\BbbR^r$ dan $J\subseteq \BbbR^s$ merupakan
interval setengah terbuka, maka
$\phi[I\times J]\subseteq\BbbR^{r+s}$ juga merupakan interval setengah terbuka, dan

\Centerline{$\mu_{r+s}(\phi[I\times J])=\mu_rI\cdot\mu_sJ$;}

\noindent ini segera diperoleh dari definisi ukuran Lebesgue
suatu interval. (Di sini saya berbicara tentang interval `setengah terbuka', yaitu
interval berbentuk
$\prod_{1\le j\le r}\coint{\alpha_j,\beta_j}$, karena saya menggunakannya dalam
definisi ukuran Lebesgue dalam \S115. Jika Anda lebih suka
bekerja dengan interval terbuka atau interval tertutup, tidak ada perbedaan.)
Perhatikan pula bahwa setiap interval setengah terbuka dalam $\BbbR^{r+s}$ dapat dinyatakan
sebagai $\phi[I\times J]$ untuk $I$, $J$ yang sesuai.

\medskip

{\bf (b)} Untuk sebarang $A\subseteq\BbbR^{r+s}$,
$\theta(\phi^{-1}[A])\le\mu^*_{r+s}(A)$. \Prf\ Untuk sebarang $\epsilon>0$,
terdapat barisan $\sequencen{K_n}$ dari interval setengah terbuka dalam $\Bbb
R^{r+s}$ sedemikian sehingga $A\subseteq\bigcup_{n\in\Bbb N}K_n$ dan
$\sum_{n=0}^{\infty}\mu_{r+s}(K_n)\le\mu^*_{r+s}(A)+\epsilon$.
Nyatakan setiap $K_n$ sebagai $\phi[I_n\times J_n]$, dengan $I_n$ dan $J_n$
interval setengah terbuka masing-masing dalam $\BbbR^r$ dan $\BbbR^s$; maka
$\phi^{-1}[A]\subseteq\bigcup_{n\in\Bbb N}I_n\times J_n$, sehingga

\Centerline{$\theta(\phi^{-1}[A])
\le\sum_{n=0}^{\infty}\mu_rI_n\cdot\mu_sJ_n
=\sum_{n=0}^{\infty}\mu_{r+s}(K_n)\le\mu^*_{r+s}(A)+\epsilon$.}

\noindent Karena $\epsilon$ sembarang, kita memperoleh hasilnya.\ \Qed

\medskip

{\bf (c)} Jika $E\subseteq\BbbR^r$ dan $F\subseteq\BbbR^s$ terukur,
maka $\mu^*_{r+s}(\phi[E\times F])\le\mu_rE\cdot\mu_sF$.

\medskip

\Prf\ {\bf (i)} Pertama-tama tinjau kasus $\mu_rE<\infty$, $\mu_sF<\infty$.
Dalam hal ini, untuk $\epsilon>0$ yang diberikan, terdapat barisan $\sequencen{I_n}$,
$\sequencen{J_n}$ dari interval setengah terbuka sedemikian sehingga
$E\subseteq\bigcup_{n\in\Bbb N}I_n$,
$F\subseteq\bigcup_{n\in\Bbb N}J_n$,

\Centerline{$\sum_{n=0}^{\infty}\mu_rI_n\le\mu^*_rE+\epsilon
=\mu_rE+\epsilon$,}

\Centerline{$\sum_{n=0}^{\infty}\mu_sJ_n\le\mu^*_sF+\epsilon=\mu_sF
+\epsilon$.}

\noindent Dengan demikian,
$E\times F\subseteq\bigcup_{m,n\in\Bbb N}I_m\times J_n$ dan
$\phi[E\times F]\subseteq\bigcup_{m,n\in\Bbb N}\phi[I_m\times J_n]$,
sehingga

$$\eqalign{\mu^*_{r+s}(\phi[E\times F])
&\le\sum_{m,n=0}^{\infty}\mu_{r+s}(\phi[I_m\times J_n])
=\sum_{m,n=0}^{\infty}\mu_rI_m\cdot\mu_sJ_n\cr
&=\sum_{m=0}^{\infty}\mu_rI_m\cdot\sum_{n=0}^{\infty}\mu_sJ_n
\le(\mu_rE+\epsilon)(\mu_sF+\epsilon).\cr}$$

\noindent Karena $\epsilon$ sembarang,
kita memperoleh hasilnya.

\medskip

\quad{\bf (ii)}
Berikutnya, jika $\mu_rE=0$, terdapat barisan $\sequencen{F_n}$ dari himpunan
berukuran hingga yang menutupi $\BbbR^s\supseteq F$, sehingga

\Centerline{$\mu^*_{r+s}(\phi[E\times F])
\le\sum_{n=0}^{\infty}\mu^*_{r+s}(\phi[E\times F_n])
\le\sum_{n=0}^{\infty}\mu_rE\cdot\mu_sF_n
=0
=\mu_rE\cdot\mu_sF$.}

\noindent Demikian pula,
$\mu^*_{r+s}(\phi[E\times F])\le\mu_rE\cdot\mu_sF$ jika $\mu_sF=0$.
Satu-satunya kasus yang tersisa adalah kasus ketika
$\mu_rE$, $\mu_sF$ keduanya bernilai positif dan salah satunya
tak hingga; tetapi dalam hal ini $\mu_rE\cdot\mu_sF=\infty$, sehingga jelas
$\mu^*_{r+s}(\phi[E\times F])\le\mu_rE\cdot\mu_sF$.\ \Qed

\medskip

{\bf (d)} Jika $A\subseteq\BbbR^{r+s}$, maka
$\mu^*_{r+s}(A)\le\theta(\phi^{-1}[A])$. \Prf\ Untuk $\epsilon>0$ yang diberikan,
terdapat barisan $\sequencen{E_n}$, $\sequencen{F_n}$ dari himpunan terukur
masing-masing dalam $\BbbR^r$, $\BbbR^s$ sedemikian sehingga
$\phi^{-1}[A]\subseteq\bigcup_{n\in\Bbb N}E_n\times F_n$ dan
$\sum_{n=0}^{\infty}\mu_rE_n\cdot\mu_sF_n
\le\theta(\phi^{-1}[A])+\epsilon$.
Sekarang $A\subseteq\bigcup_{n\in\Bbb N}\phi[E_n\times F_n]$, sehingga

\Centerline{$\mu^*_{r+s}(A)
\le\sum_{n=0}^{\infty}\mu^*_{r+s}(\phi[E_n\times F_n])
\le\sum_{n=0}^{\infty}\mu_rE_n\cdot\mu_sF_n
\le\theta(\phi^{-1}[A])+\epsilon$.}

\noindent Karena $\epsilon$ sembarang, kita memperoleh hasilnya.\ \Qed

\medskip

{\bf (e)} Dengan menggabungkan (b) dan (d), kita memperoleh
$\theta(\phi^{-1}[A])=\mu^*_{r+s}(A)$ untuk setiap $A\subseteq\Bbb
R^{r+s}$. Jadi $\theta$ pada $\BbbR^r\times\BbbR^s$ berkorespondensi
tepat dengan $\mu^*_{r+s}$ pada $\BbbR^{r+s}$. Maka ukuran-ukuran terkait
$\lambda_0$, $\mu_{r+s}$ harus berkorespondensi dengan cara yang sama,
dengan menuliskan $\lambda_0$ untuk ukuran produk primitif. Tetapi 251K memberi tahu
kita bahwa $\lambda_0=\lambda$, sehingga kita memperoleh hasilnya.
}%end of proof of 251N

\vleader{48pt}{251O}{}\cmmnt{ Sebenarnya, sebagian besar
penerapan
konstruksi di sini adalah pada subruang dari ruang Euklides, bukan
pada keseluruhan produk $\BbbR^r\times\BbbR^s$. Tidaklah terlalu
sulit untuk menuliskan 251N agar menangani
subruang sebarang, tetapi saya lebih memilih memberikan uraian yang lebih umum mengenai
produk ukuran subruang, karena menurut saya uraian ini memperjelas
metodenya. Saya mulai dengan hasil langsung tentang ruang yang dapat dilokalkan secara
ketat.

\medskip

\noindent}{\bf Proposisi} Misalkan $(X,\Sigma,\mu)$ dan $(Y,\Tau,\nu)$
ruang ukur yang dapat dilokalkan secara ketat. Maka ukuran produk c.l.d.\ pada
$X\times Y$ dapat dilokalkan secara ketat; selain itu, jika
$\langle X_i\rangle_{i\in I}$ dan $\langle Y_j\rangle_{j\in J}$ merupakan
dekomposisi dari $X$ dan $Y$ masing-masing,
$\langle X_i\times Y_j\rangle_{(i,j)\in I\times J}$ merupakan dekomposisi dari
$X\times Y$.

\proof{ Misalkan $\langle X_i\rangle_{i\in I}$ dan
$\langle Y_j\rangle_{j\in J}$ dekomposisi dari $X$, $Y$ masing-masing.
Maka $\langle X_i\times Y_j\rangle_{(i,j)\in I\times J}$ merupakan partisi dari
$X\times Y$ menjadi himpunan-himpunan terukur yang berukuran hingga. Jika
$W\subseteq X\times Y$ dan $\lambda W>0$, terdapat himpunan $E\in\Sigma$,
$F\in\Tau$ sedemikian sehingga $\mu E<\infty$, $\nu F<\infty$ dan
$\lambda(W\cap(E\times F))>0$. Kita mengetahui bahwa
$\mu E=\sum_{i\in I}\mu(E\cap X_i)$ dan
$\nu F=\sum_{j\in J}\nu(F\cap Y_j)$, sehingga harus terdapat himpunan berhingga
$I_0\subseteq I$, $J_0\subseteq J$ sedemikian sehingga

\Centerline{$\mu E\cdot\nu F
  -(\sum_{i\in I_0}\mu(E\cap X_i))(\sum_{j\in J_0}\nu(F\cap Y_j))
<\lambda(W\cap(E\times F))$.}

\noindent Dengan menetapkan $E'=\bigcup_{i\in I_0}X_i$ dan $F'=\bigcup_{j\in J_0}Y_j$
kita mempunyai

\Centerline{$\lambda((E\times F)\setminus(E'\times F'))
=\lambda(E\times F)-\lambda((E\cap E')\times(F\cap
F'))<\lambda(W\cap(E\times F))$,}

\noindent sehingga $\lambda(W\cap(E'\times F'))>0$. Oleh karena itu harus
terdapat $i\in I_0$, $j\in J_0$ sedemikian sehingga
$\lambda(W\cap(X_i\times Y_j))>0$.

Ini menunjukkan bahwa $\{X_i\times Y_j:i\in I,\,j\in J\}$ memenuhi
kriteria 213O, sehingga $\lambda$, karena lengkap dan ditentukan secara
lokal, harus dapat dilokalkan secara ketat. Karena
$\langle X_i\times Y_j\rangle_{(i,j)\in I\times J}$ menutupi $X\times Y$,
keluarga ini benar-benar merupakan dekomposisi dari $X\times Y$ (213Ob).
}

\leader{251P}{Lema} Misalkan $(X,\Sigma,\mu)$ dan $(Y,\Tau,\nu)$ ruang
ukur, dan $\lambda$ ukuran produk c.l.d.\ pada
$X\times Y$. Misalkan $\lambda^*$ merupakan
ukuran luar yang bersesuaian\cmmnt{ (132B)}. Maka

\Centerline{$\lambda^*C=\sup\{\theta(C\cap(E\times
F)):E\in\Sigma,\,F\in\Tau,\,\mu E<\infty,\,\nu F<\infty\}$}

\noindent untuk setiap
$C\subseteq X\times Y$, dengan $\theta$ ukuran luar dari 251A.

\proof{ Tuliskan $\Lambda$ untuk domain $\lambda$, $\Sigma^f$ untuk
$\{E:E\in\Sigma,\,\mu E<\infty\}$, $\Tau^f$ untuk $\{F:F\in\Tau,\,\nu
F<\infty\}$; tetapkan $u=\sup\{\theta(C\cap(E\times F)):E\in\Sigma^f$,
$F\in\Tau^f\}$.

\medskip

{\bf (a)} Jika $C\subseteq W\in\Lambda$, $E\in\Sigma^f$ dan $F\in\Tau^f$,
maka

$$\eqalignno{\theta(C\cap(E\times F))
&\le\theta(W\cap(E\times F))
=\lambda_0(W\cap(E\times F))\cr
\noalign{\noindent (dengan $\lambda_0$ ukuran produk primitif)}
&\le\lambda W.\cr}$$

\noindent Karena $E$ dan $F$ sembarang, $u\le\lambda W$; karena $W$
sembarang, $u\le\lambda^*C$.

\medskip

{\bf (b)} Jika $u=\infty$, maka tentu saja $\lambda^*C=u$. Jika tidak, misalkan
$\sequencen{E_n}$, $\sequencen{F_n}$ merupakan
barisan masing-masing dalam $\Sigma^f$, $\Tau^f$ sedemikian sehingga

\Centerline{$u=\sup_{n\in\Bbb N}\theta(C\cap(E_n\times F_n))$.}

\noindent Tinjau $C'=C\setminus(\bigcup_{n\in\Bbb
N}E_n\times\bigcup_{n\in\Bbb
N}F_n)$. Jika $E\in\Sigma^f$ dan $F\in\Tau^f$, maka untuk setiap $n\in\Bbb
N$ kita mempunyai

$$\eqalignno{u
&\ge\theta(C\cap((E\cup E_n)\times(F\cup F_n)))\cr
&=\theta(C\cap((E\cup E_n)\times(F\cup F_n))\cap(E_n\times F_n))\cr
&{\hskip 10em}+\theta(C\cap((E\cup E_n)\times(F\cup F_n))
  \setminus(E_n\times F_n))\cr
\noalign{\noindent (karena $E_n\times F_n\in\Lambda$, menurut 251E)}
&\ge\theta(C\cap(E_n\times F_n))
  +\theta(C'\cap(E\times F)).\cr}$$

\noindent Dengan mengambil supremum dari ungkapan di ruas kanan ketika $n$
bervariasi, kita memperoleh $u\ge u+\theta(C'\cap(E\times F))$, sehingga

\Centerline{$\lambda(C'\cap(E\times F))=\theta(C'\cap(E\times F))=0$.}

\noindent Karena $E$ dan $F$ sembarang, $\lambda C'=0$.

Tetapi ini berarti bahwa

$$\eqalignno{\lambda^*C
&\le\lambda^*
  (C\cap(\bigcup_{n\in\Bbb N}E_n\times\bigcup_{n\in\Bbb N}F_n))
  +\lambda^*C'\cr
&=\lim_{n\to\infty}
  \lambda^*(C\cap(\bigcup_{i\le n}E_i\times\bigcup_{i\le n}F_i))\cr
\noalign{\noindent (menggunakan 132Ae)}
&\le u,\cr}$$

\noindent sebagaimana diperlukan.
}%end of proof of 251P

\leader{251Q}{\bf Proposisi} Misalkan $(X,\Sigma,\mu)$ dan $(Y,\Tau,\nu)$
ruang ukur, serta $A\subseteq X$, $B\subseteq Y$ himpunan bagian; tuliskan
$\mu_A$, $\nu_B$ untuk ukuran subruang masing-masing pada $A$, $B$.
Misalkan $\lambda$ ukuran produk c.l.d.\ pada $X\times Y$, dan
$\lambda^{\#}$
ukuran subruang yang diinduksinya pada $A\times B$. Misalkan $\tilde\lambda$
ukuran produk c.l.d.\ dari $\mu_A$ dan $\nu_B$ pada $A\times B$. Maka

(i) $\tilde\lambda$ memperluas $\lambda^{\#}$.

(ii) Jika

\inset{{\it salah satu} ($\alpha$) $A\in\Sigma$ dan $B\in\Tau$}

\inset{{\it atau} ($\beta$) $A$ dan $B$ keduanya dapat ditutupi oleh barisan
himpunan berukuran hingga}

\inset{{\it atau} ($\gamma$) $\mu$ dan $\nu$ keduanya dapat dilokalkan secara
ketat,}

\noindent maka $\tilde\lambda=\lambda^{\#}$.

\proof{ Misalkan $\theta$ ukuran luar pada $X\times Y$ yang didefinisikan dari
$\mu$ dan $\nu$ melalui rumus dalam 251A, dan $\tilde\theta$ ukuran luar
pada $A\times B$ yang didefinisikan secara serupa dari $\mu_A$ dan $\nu_B$.
Tuliskan $\Lambda$ untuk domain dari $\lambda$,
$\tilde\Lambda$ untuk domain dari $\tilde\lambda$, dan
$\Lambda^{\#}=\{W\cap(A\times B):W\in\Lambda\}$ untuk
domain dari $\lambda^{\#}$.
Tetapkan $\Sigma^f=\{E:\mu E<\infty\}$, $\Tau^f=\{F:\nu F<\infty\}$.

\medskip

{\bf (a)} Hal pertama yang perlu diperhatikan adalah bahwa $\tilde\theta C=\theta C$
untuk setiap $C\subseteq A\times B$. \Prf\ (i) Jika $\sequencen{E_n}$
dan $\sequencen{F_n}$ merupakan barisan masing-masing dalam $\Sigma$, $\Tau$
sedemikian sehingga $C\subseteq\bigcup_{n\in\Bbb N}E_n\times F_n$, maka

\Centerline{$C=C\cap(A\times B)\subseteq\bigcup_{n\in\Bbb N}(E_n\cap
A)\times(F_n\cap B)$,}

\noindent sehingga

$$\eqalign{\tilde\theta C
&\le\sum_{n=0}^{\infty}\mu_A(E_n\cap A)\cdot\nu_B(F_n\cap B)\cr
&=\sum_{n=0}^{\infty}\mu^*(E_n\cap A)\cdot\nu^*(F_n\cap B)
\le\sum_{n=0}^{\infty}\mu E_n\cdot\nu F_n.\cr}$$

\noindent Karena $\sequencen{E_n}$ dan $\sequencen{F_n}$ sembarang,
$\tilde\theta C\le\theta C$. (ii) Jika $\sequencen{\tilde E_n}$,
$\sequencen{\tilde F_n}$ merupakan barisan masing-masing dalam $\Sigma_A=\dom\mu_A$,
$\Tau_B=\dom\nu_B$ sedemikian sehingga
$C\subseteq\bigcup_{n\in\Bbb N}\tilde E_n\times\tilde F_n$, maka untuk
setiap $n\in\Bbb N$ kita dapat memilih
$E_n\in\Sigma$, $F_n\in\Tau$ sedemikian sehingga

\Centerline{$\tilde E_n\subseteq E_n$,\quad $\mu E_n
=\mu^*\tilde E_n=\mu_A\tilde E_n$,}

\Centerline{$\tilde F_n\subseteq F_n$,
\quad $\nu F_n=\nu^*\tilde F_n=\nu_B\tilde F_n$,}

\noindent dan sekarang

\Centerline{$\theta C\le\sum_{n=0}^{\infty}\mu E_n\cdot\nu F_n
=\sum_{n=0}^{\infty}\mu_A\tilde E_n\cdot\nu_B\tilde F_n$.}

\noindent Karena $\sequencen{\tilde E_n}$, $\sequencen{\tilde F_n}$
sembarang, $\theta C\le\tilde\theta C$.\ \Qed

\medskip

{\bf (b)} Karena itu, $\Lambda^{\#}\subseteq\tilde\Lambda$. \Prf\
Misalkan $V\in\Lambda^{\#}$ dan $C\subseteq A\times B$.
Dalam
hal ini terdapat $W\in\Lambda$ sedemikian sehingga $V=W\cap(A\times B)$. Jadi

\Centerline{$\tilde\theta(C\cap V)+\tilde\theta(C\setminus V)
=\theta(C\cap W)+\theta(C\setminus W)
=\theta C
=\tilde\theta C$.}

\noindent Karena $C$ sembarang, $V\in\tilde\Lambda$.\ \Qed

Dengan demikian, untuk $V\in\Lambda^{\#}$,

$$\eqalignno{\lambda^{\#}V
&=\lambda^*V
=\sup\{\theta(V\cap(E\times F)):E\in\Sigma^f,\,F\in\Tau^f\}\cr
&=\sup\{\theta(V\cap(\tilde E\times\tilde F)):\tilde
E\in\Sigma_A,\,\tilde F\in\Tau_B,\,\mu_A\tilde E<\infty,\,\nu_B\tilde
F<\infty\}\cr
&=\sup\{\tilde\theta(V\cap(\tilde E\times\tilde F)):\tilde
E\in\Sigma_A,\,\tilde F\in\Tau_B,\,\mu_A\tilde E<\infty,\,\nu_B\tilde
F<\infty\}
=\tilde\lambda V,\cr}$$

\noindent dengan menggunakan 251P dua kali.

Ini membuktikan bagian (i) dari proposisi.

\medskip

{\bf (c)} Hal berikutnya yang perlu diperiksa adalah bahwa jika $V\in\tilde\Lambda$ dan
$V\subseteq E\times F$ dengan $E\in\Sigma^f$ dan $F\in\Tau^f$, maka
$V\in\Lambda^{\#}$. \Prf\ Misalkan $W\subseteq E\times F$ merupakan selubung
terukur dari $V$ terhadap $\lambda$ (132Ee). Maka

$$\eqalignno{\theta(W\cap(A\times B)\setminus V)
&=\tilde\theta(W\cap(A\times B)\setminus V)
=\tilde\lambda(W\cap(A\times B)\setminus V)\cr
\noalign{\noindent (karena $W\cap(A\times
B)\in\Lambda^{\#}\subseteq\tilde\Lambda$, $V\in\tilde\Lambda$)}
&=\tilde\lambda(W\cap(A\times B))-\tilde\lambda V
=\tilde\theta(W\cap(A\times B))-\tilde\theta V\cr
&=\theta(W\cap(A\times B))-\theta V
=\lambda^*(W\cap(A\times B))-\lambda^* V\cr
&\le\lambda W-\lambda^* V
=0.\cr}$$

\noindent Tetapi ini berarti bahwa $W'=W\cap(A\times B)\setminus V\in\Lambda$
dan $V=(A\times B)\cap(W\setminus W')$ termasuk dalam $\Lambda^{\#}$.\ \Qed

\medskip

{\bf (d)} Sekarang tetapkan sebarang $V\in\tilde\Lambda$, dan periksa syarat
($\alpha$)-($\gamma$) dari bagian (ii) proposisi.

\medskip

\quad\grheada\ Jika $A\in\Sigma$ dan $B\in\Tau$, serta $C\subseteq X\times Y$,
maka $A\times B\in\Lambda$ (251E), sehingga

$$\eqalignno{\theta(C\cap V)+\theta(C\setminus V)
&=\theta(C\cap V)+\theta((C\setminus V)\cap(A\times B))
  +\theta((C\setminus V)\setminus(A\times B))\cr
&=\tilde\theta(C\cap V)+\tilde\theta(C\cap(A\times B)\setminus V)
  +\theta(C\setminus(A\times B))\cr
&=\tilde\theta(C\cap(A\times B))
  +\theta(C\setminus(A\times B))\cr
&=\theta(C\cap(A\times B))
  +\theta(C\setminus(A\times B))
=\theta C.\cr}$$

\noindent Karena $C$ sembarang, $V\in\Lambda$, sehingga
$V=V\cap(A\times B)$ termasuk dalam $\Lambda^{\#}$.

\medskip

\quad\grheadb\ Jika $A\subseteq\bigcup_{n\in\Bbb N}E_n$ dan
$B\subseteq\bigcup_{n\in\Bbb N}F_n$ dengan semua $E_n$, $F_n$ berukuran
hingga, maka
$V=\bigcup_{m,n\in\Bbb N}V\cap(E_m\times F_n)\in\Lambda^{\#}$, menurut (c).

\medskip

\quad\grheadc\ Jika $\familyiI{X_i}$, $\family{j}{J}{Y_j}$ merupakan
dekomposisi
dari $X$, $Y$ masing-masing, maka untuk setiap $i\in I$, $j\in J$ kita mempunyai
$V\cap(X_i\times Y_j)\in\Lambda^{\#}$, yakni, terdapat
$W_{ij}\in\Lambda$ sedemikian sehingga
$V\cap(X_i\times Y_j)=W_{ij}\cap(A\times B)$.
Sekarang $\family{(i,j)}{I\times J}{X_i\times Y_j}$ merupakan dekomposisi dari
$X\times Y$ untuk $\lambda$ (251O), sehingga

\Centerline{$W=\bigcup_{i\in I,j\in J}W_{ij}\cap(X_i\times Y_j)
\in\Lambda$,}

\noindent dan $V=W\cap(A\times B)\in\Lambda^{\#}$.

\medskip

{\bf (e)} Jadi masing-masing dari ketiga syarat itu cukup untuk memastikan bahwa
$\tilde\Lambda=\Lambda^{\#}$, dan dalam hal ini (a) memberi tahu kita bahwa
$\tilde\lambda=\lambda^{\#}$.
}%end of proof of 251Q
\leader{251R}{Akibat} Misalkan $r$, $s\ge 1$ bilangan bulat, dan
$\phi:\Bbb R^r\times\BbbR^s\to\BbbR^{r+s}$ bijeksi alamiah. Jika
$A\subseteq\BbbR^r$ dan $B\subseteq\BbbR^s$, maka pembatasan
$\phi$ pada $A\times B$ mengidentifikasikan produk ukuran Lebesgue pada $A$
dan ukuran Lebesgue pada $B$ dengan ukuran Lebesgue pada $\phi[A\times
B]\subseteq\BbbR^{r+s}$.

\cmmnt{\medskip

\noindent{\bf Catatan} Perhatikan bahwa dengan `ukuran Lebesgue pada $A$' yang saya maksud
adalah ukuran subruang $\mu_{rA}$ pada $A$ yang diinduksi oleh ukuran Lebesgue
berdimensi $r$, $\mu_r$, pada $\BbbR^r$, terlepas dari apakah $A$ sendiri
merupakan himpunan terukur atau bukan.
}%end of comment

\proof{ Menurut 251Q, dengan menggunakan salah satu syarat
(ii-$\beta$) atau (ii-$\gamma$), ukuran produk $\tilde\lambda$ pada
$A\times B$ tidak lain adalah ukuran subruang $\lambda^{\#}$ pada $A\times B$
yang diinduksi oleh
ukuran produk $\lambda$ pada $\BbbR^r\times\BbbR^s$. Namun, menurut 251N kita
tahu bahwa $\phi$ merupakan isomorfisme antara $(\BbbR^r\times\Bbb
R^s,\lambda)$ dan $(\BbbR^{r+s},\mu_{r+s})$; jadi ia juga harus
mengidentifikasikan $\tilde\lambda$ dengan ukuran subruang pada $\phi[A\times B]$.
}

\leader{251S}{Akibat} Misalkan $(X,\Sigma,\mu)$ dan $(Y,\Tau,\nu)$
ruang ukur, dan $\lambda$ ukuran produk c.l.d.\
pada $X\times Y$. Jika
$A\subseteq X$ dan $B\subseteq Y$ dapat ditutupi oleh barisan-barisan himpunan
berukuran hingga, maka $\lambda^*(A\times
B)=\mu^*A\cdot\nu^*B$.

\proof{ Dalam bahasa 251Q,

$$\eqalignno{\lambda^*(A\times B)
&=\lambda^{\#}(A\times B)
=\tilde\lambda(A\times B)
=\mu_AA\cdot\nu_BB\cr
\noalign{\noindent (menurut 251K dan 251E)}
&=\mu^*A\cdot\nu^*B.\cr}$$
}%end of proof of 251S

\leader{251T}{}\cmmnt{ Proposisi berikut memberikan gambaran tentang bagaimana
definisi-definisi teknis di sini saling berpadu.

\medskip

\noindent}{\bf Proposisi} Misalkan $(X,\Sigma,\mu)$ dan $(Y,\Tau,\nu)$
ruang ukur. Tuliskan $(X,\hat\Sigma,\hat\mu)$ dan
$(X,\tilde\Sigma,\tilde\mu)$ untuk pelengkapan dan versi c.l.d.\ dari
$(X,\Sigma,\mu)$\cmmnt{ (212C, 213E)}. Misalkan $\lambda$, $\hat\lambda$
dan $\tilde\lambda$ merupakan tiga ukuran produk c.l.d.\ pada $X\times Y$
yang diperoleh dari ketiga pasangan ukuran faktor $(\mu,\nu)$,
$(\hat\mu,\nu)$ dan $(\tilde\mu,\nu)$. Maka
$\lambda=\hat\lambda=\tilde\lambda$.

\proof{ Tuliskan $\Lambda$, $\hat\Lambda$ dan $\tilde\Lambda$ untuk
domain
$\lambda$, $\hat\lambda$, $\tilde\lambda$ secara berturut-turut; dan
$\theta$,
$\hat\theta$, $\tilde\theta$ untuk ukuran-ukuran luar pada $X\times Y$
yang diperoleh
dengan rumus 251A dari ketiga pasangan ukuran faktor tersebut.

\medskip

{\bf (a)} Jika $E\in\Sigma$ dan $\mu E<\infty$, maka $\theta$,
$\hat\theta$
dan $\tilde\theta$ bertepatan pada subhimpunan-subhimpunan $E\times Y$. \Prf\ Ambil
$A\subseteq E\times Y$ dan $\epsilon>0$.

\medskip

\quad{\bf (i)} Terdapat barisan $\sequencen{E_n}$ dalam $\Sigma$,
$\sequencen{F_n}$ dalam $\Tau$ sedemikian sehingga
$A\subseteq\bigcup_{n\in\Bbb N}E_n\times F_n$ dan
$\sum_{n=0}^{\infty}\mu E_n\cdot\nu F_n\le\theta A+\epsilon$. Sekarang
$\tilde\mu E_n\le\mu E_n$ untuk setiap $n$ (213Fb), sehingga

\Centerline{$\tilde\theta A
\le\sum_{n=0}^{\infty}\tilde\mu E_n\cdot\nu F_n
\le\sum_{n=0}^{\infty}\mu E_n\cdot\nu F_n\le\theta A+\epsilon$.}

\medskip

\quad{\bf (ii)} Terdapat barisan $\sequencen{\hat E_n}$ dalam
$\hat\Sigma$, $\sequencen{\hat F_n}$ dalam $\Tau$ sedemikian sehingga
$A\subseteq\bigcup_{n\in\Bbb N}\hat E_n\times\hat F_n$ dan
$\sum_{n=0}^{\infty}\hat\mu\hat E_n\cdot\nu\hat F_n
\le\hat\theta A+\epsilon$. Sekarang, untuk setiap $n$ terdapat
$E'_n\in\Sigma$ sedemikian sehingga $\hat E_n\subseteq E'_n$ dan
$\mu E'_n=\hat\mu\hat E_n$, sehingga

\Centerline{$\theta A\le\sum_{n=0}^{\infty}\mu E'_n\cdot\nu\hat F_n
=\sum_{n=0}^{\infty}\hat\mu\hat E_n\cdot\nu\hat F_n
\le\hat\theta A+\epsilon$.}

\medskip

\quad{\bf (iii)} Terdapat barisan $\sequencen{\tilde E_n}$ dalam
$\tilde\Sigma$, $\sequencen{\tilde F_n}$ dalam $\Tau$ sedemikian sehingga
$A\subseteq\bigcup_{n\in\Bbb N}\tilde E_n\times\tilde F_n$ dan
$\sum_{n=0}^{\infty}\tilde\mu\tilde E_n\cdot\nu\tilde F_n
\le\tilde\theta A+\epsilon$. Sekarang, untuk setiap $n$,
$\tilde E_n\cap E\in\hat\Sigma$, sehingga

\Centerline{$\hat\theta A
\le\sum_{n=0}^{\infty}\hat\mu(\tilde E_n\cap E)\cdot\nu\tilde F_n
\le\sum_{n=0}^{\infty}\tilde\mu\tilde E_n\cdot\nu\tilde F_n
\le\tilde\theta A+\epsilon$.}

\medskip

\quad{\bf (iv)} Karena $A$ dan $\epsilon$ sebarang,
$\theta=\hat\theta=\tilde\theta$ pada $\Cal P(E\times Y)$.\ \Qed

\medskip

{\bf (b)} Akibatnya, ukuran-ukuran luar $\lambda^*$, $\hat\lambda^*$
dan $\tilde\lambda^*$ identik. \Prf\ Gunakan 251P. Ambil
$A\subseteq X\times Y$, $E\in\Sigma$, $\hat E\in\hat\Sigma$,
$\tilde E\in\tilde\Sigma$,
$F\in\Tau$ sedemikian sehingga $\mu E$, $\hat\mu\hat E$, $\tilde\mu\tilde E$ dan
$\nu F$ semuanya hingga. Maka

\medskip

\quad{\bf (i)}

\Centerline{$\theta(A\cap(E\times F))=\hat\theta(A\cap(E\times F))
\le\hat\lambda^*A$,
\quad$\theta(A\cap(E\times F))=\tilde\theta(A\cap(E\times F))
\le\tilde\lambda^*A$}

\noindent karena $\hat\mu E$ dan $\tilde\mu E$ keduanya hingga.

\medskip

\quad{\bf (ii)} Terdapat $E'\in\Sigma$ sedemikian sehingga $\hat E\subseteq E'$
dan $\mu E'<\infty$, sehingga

\Centerline{$\hat\theta(A\cap(\hat E\times F))
\le\hat\theta(A\cap(E'\times F))
=\theta(A\cap(E'\times F))
\le\lambda^*A$.}

\medskip

\quad{\bf (iii)} Terdapat $E''\in\Sigma$ sedemikian sehingga
$E''\subseteq\tilde E$ dan $\tilde\mu(\tilde E\setminus E'')=0$ (213Fc),
sehingga
$\tilde\theta((\tilde E\setminus E'')\times Y)=0$ dan $\mu E''<\infty$;
dengan demikian

\Centerline{$\tilde\theta(A\cap(\tilde E\times F))
=\tilde\theta(A\cap(E''\times F))
=\theta(A\cap(E''\times F))
\le\lambda^*A$.}

\medskip

\quad{\bf (iv)} Dengan mengambil supremum atas $E$, $\hat E$, $\tilde E$ dan
$F$,
kita memperoleh

\Centerline{$\lambda^*A\le\hat\lambda^*A$,
\quad$\lambda^*A\le\tilde\lambda^*A$,
\quad$\hat\lambda^*A\le\lambda^*A$,
\quad$\tilde\lambda^*A\le\lambda^*A$.}

\noindent Karena $A$ sebarang,
$\lambda^*=\hat\lambda^*=\tilde\lambda^*$.\ \Qed

\medskip

{\bf (c)} Sekarang $\lambda$, $\hat\lambda$ dan $\tilde\lambda$ semuanya
lengkap dan ditentukan secara lokal, sehingga menurut 213C ketiganya merupakan ukuran
yang didefinisikan dengan metode
\Caratheodory\ dari ukuran luar masing-masing, dan karena itu identik.
}%end of proof of 251T

\leader{251U}{}\cmmnt{ `Jelas', dan merupakan akibat mudah dari
teorema-teorema yang telah dibuktikan sejauh ini, bahwa himpunan $\{(x,x):x\in\Bbb R\}$
terabaikan terhadap ukuran Lebesgue pada $\BbbR^2$. Hasil yang bersesuaian
berlaku pada kuadrat sebarang ruang ukur {\it tanpa atom}.

\medskip

\noindent}{\bf Proposisi} Misalkan $(X,\Sigma,\mu)$ ruang ukur tanpa atom,
dan misalkan $\lambda$ ukuran produk c.l.d.\ pada $X\times X$. Maka
$\Delta=\{(x,x):x\in X\}$ terabaikan terhadap $\lambda$.

\proof{ Misalkan $E$, $F\in\Sigma$ himpunan-himpunan berukuran hingga, dan
$n\in\Bbb N$. Dengan menerapkan 215D berulang kali, kita dapat memperoleh keluarga
saling lepas $\ofamily{i}{n}{F_i}$ dari subhimpunan-subhimpunan terukur dari $F$
sedemikian sehingga $\mu F_i=\Bover{\mu F}{n+1}$ untuk setiap $i$; dengan menetapkan
$F_n=F\setminus\bigcup_{i<n}F_i$, kita juga mempunyai
$\mu F_n=\Bover{\mu F}{n+1}$. Sekarang

\Centerline{$\Delta\cap(E\times F)
\subseteq\bigcup_{i\le n}(E\cap F_i)\times F_i$,}

\noindent sehingga

\Centerline{$\lambda^*(\Delta\cap(E\times F))
\le\sum_{i=0}^n\mu(E\cap F_i)\cdot\mu F_i
=\Bover{\mu F}{n+1}\sum_{i=0}^n\mu(E\cap F_i)
\le\Bover1{n+1}\mu E\cdot\mu F$.}

\noindent Karena $n$ sebarang, $\lambda(\Delta\cap(E\times F))=0$; karena
$E$ dan $F$ sebarang, $\lambda\Delta=0$.
}%end of proof of 251U

%\mu\times(\nu_1+\nu_2)=(\mu\times\nu_1)+(\mu\times\nu_2)
%may need restrictions?

\leader{*251W}{Produk lebih dari dua ruang}\cmmnt{ Seluruh
bagian ini dapat diulangi untuk produk berhingga sebarang. Pekerjaannya
cukup banyak, tetapi tidak diperlukan gagasan baru. Pada saat kita memerlukan
konstruksi umum ini secara formal, konstruksi itu semestinya sudah terasa sangat wajar,
dan saya kira halaman berikut tidak perlu dikerjakan sebelum melanjutkan,
terutama karena produk ruang-ruang {\it probabilitas}
dibahas dalam \S254. Namun, demi kelengkapan, dan untuk membantu
menemukan hasil ketika penerapannya muncul, saya mencantumkannya di sini. Hasil-hasil ini
tentu saja membentuk seperangkat latihan yang sangat mendidik. Bagian-bagian terpenting
diulangi dalam 251Xe-251Xf.

} Misalkan $\familyiI{(X_i,\Sigma_i,\mu_i)}$ keluarga berhingga ruang ukur,
dan tetapkan $X=\prod_{i\in I}X_i$. Tuliskan
$\Sigma_i^f=\{E:E\in\Sigma_i,\,\mu_iE<\infty\}$ untuk setiap $i\in I$.

\spheader 251Wa Untuk $A\subseteq X$ tetapkan

$$\theta A=\inf\{\sum_{n=0}^{\infty}\prod_{i\in I}\mu_iE_{ni}:
E_{ni}\in\Sigma_i\Forall i\in I,\,n\in\Bbb N,\,
A\subseteq\bigcup_{n\in\Bbb N}\prod_{i\in I}E_{ni}\}.$$

\noindent Maka $\theta$ merupakan ukuran luar pada $X$. Misalkan $\lambda_0$
ukuran pada $X$ yang diturunkan dengan metode \Caratheodory\ dari $\theta$, dan
$\Lambda$ domainnya.

\spheader 251Wb Jika $\familyiI{X_i}$ keluarga berhingga himpunan dan
$\Sigma_i$ merupakan $\sigma$-aljabar subhimpunan $X_i$ untuk setiap $i\in I$,
maka $\Tensorhat_{i\in I}\Sigma_i$ adalah $\sigma$-aljabar subhimpunan
$X=\prod_{i\in I}X_i$ yang dibangkitkan oleh $\{\prod_{i\in I}E_i:E_i\in\Sigma_i$
untuk setiap $i\in I\}$. \cmmnt{(Untuk konstruksi yang bersesuaian
ketika $I$ tak hingga, lihat 254E.)}

\spheader 251Wc $\lambda_0(\prod_{i\in I}E_i)$ terdefinisi dan sama dengan
$\prod_{i\in I}\mu_iE_i$ kapan pun $E_i\in\Sigma_i$ untuk setiap $i\in I$.

\spheader 251Wd {\bf ukuran produk c.l.d.\ } pada $X$ adalah ukuran
$\lambda$ yang didefinisikan dengan menetapkan

\Centerline{$\lambda W
=\sup\{\lambda_0(W\cap\prod_{i\in I}E_i):
  E_i\in\Sigma_i^f$ untuk setiap $i\in I\}$}

\noindent untuk $W\in\Lambda$. Jika $I$ kosong, \cmmnt{sehingga
$X=\{\emptyset\}$, maka konvensi yang sesuai ialah} menetapkan
$\lambda X=1$.

\spheader 251We Jika $H\subseteq X$, maka $H\in\Lambda$ jika dan hanya jika
$H\cap\prod_{i\in I}E_i\in\Lambda$ kapan pun $E_i\in\Sigma_i^f$ untuk setiap
$i\in I$.

\spheader 251Wf(i) $\Tensorhat_{i\in I}\Sigma_i\subseteq\Lambda$ dan
$\lambda(\prod_{i\in I}E_i)=\prod_{i\in I}\mu_iE_i$ kapan pun
$E_i\in\Sigma_i^f$ untuk setiap $i$.

\quad(ii) Untuk setiap $W\in\Lambda$ terdapat $V\in\Tensorhat_{i\in
I}\Sigma_i$ sedemikian sehingga $V\subseteq W$ dan $\lambda V=\lambda W$.

\quad(iii) $\lambda$ lengkap dan ditentukan secara lokal, serta merupakan
versi c.l.d.\ dari $\lambda_0$.

\quad(iv) Jika $W\in\Lambda$ dan $\lambda W>0$, maka terdapat
$E_i\in\Sigma_i^f$, untuk $i\in I$, sedemikian sehingga
$\lambda(W\cap\prod_{i\in I}E_i)>0$.

\quad(v) Jika $W\in\Lambda$ dan $\lambda W<\infty$, maka untuk setiap
$\epsilon>0$ terdapat $n\in\Bbb N$ dan
$E_{0i},\ldots,E_{ni}\in\Sigma_i^f$, untuk setiap $i\in I$, sedemikian sehingga
$\lambda(W\symmdiff\bigcup_{k\le n}\prod_{i\in I}E_{ki})\le\epsilon$.

\spheader 251Wg Jika setiap $\mu_i$ $\sigma$-hingga, maka $\lambda$ juga demikian, dan
$\lambda=\lambda_0$ merupakan pelengkapan dari pembatasannya pada
$\Tensorhat_{i\in I}\Sigma_i$.

\spheader 251Wh Jika $\family{j}{J}{I_j}$ sebarang partisi $I$, maka
$\lambda$ dapat diidentifikasikan dengan produk c.l.d.\ dari
$\family{j}{J}{\lambda_j}$, dengan $\lambda_j$ produk c.l.d.\ dari
$\family{i}{I_j}{\mu_i}$. \cmmnt{(Lihat argumen dalam 251N dan juga
dalam 254N di bawah.)}

\spheader 251Wi Jika $I=\{1,\ldots,n\}$ dan setiap $\mu_i$ merupakan ukuran Lebesgue
pada $\Bbb R$, maka $\lambda$ dapat diidentifikasikan dengan ukuran Lebesgue
pada $\BbbR^n$.

\spheader 251Wj Jika, untuk setiap $i\in I$, kita mempunyai dekomposisi
$\family{j}{J_i}{X_{ij}}$ dari $X_i$, maka
$\langle\prod_{i\in I}X_{i,f(i)}\rangle_{f\in\prod_{i\in I}J_i}$ merupakan
dekomposisi $X$.

\spheader 251Wk Untuk sebarang $C\subseteq X$,

\Centerline{$\lambda^*C
=\sup\{\theta(C\cap\prod_{i\in I}E_i):
  E_i\in\Sigma_i^f$ untuk setiap $i\in I\}$.}

\spheader 251Wl Misalkan $A_i\subseteq X_i$ untuk setiap $i\in I$.
Tuliskan $\lambda^{\#}$ untuk ukuran subruang pada $A=\prod_{i\in I}A_i$,
dan $\tilde\lambda$ untuk produk c.l.d.\ dari ukuran-ukuran subruang pada
$A_i$. Maka $\tilde\lambda$ memperluas $\lambda^{\#}$, dan jika

\quad{\it salah satu} $A_i\in\Sigma_i$ untuk setiap $i$

\quad{\it atau} setiap $A_i$ dapat ditutupi oleh suatu barisan himpunan berukuran
hingga

\quad{\it atau} setiap $\mu_i$ dapat dilokalkan secara ketat,

\noindent maka $\tilde\lambda=\lambda^{\#}$.

\spheader 251Wm Jika $A_i\subseteq X_i$ dapat ditutupi oleh suatu barisan
himpunan berukuran hingga untuk setiap $i\in I$, maka
$\lambda^*(\prod_{i\in I}A_i)=\prod_{i\in I}\mu_i^*A_i$.

\spheader 251Wn Dengan menuliskan $\hat\mu_i$, $\tilde\mu_i$ untuk pelengkapan
dan versi c.l.d.\ dari setiap $\mu_i$, $\lambda$ merupakan produk c.l.d.\ dari
$\familyiI{\hat\mu_i}$ dan juga dari $\familyiI{\tilde\mu_i}$.

\spheader 251Wo Jika semua $(X_i,\Sigma_i,\mu_i)$ merupakan ruang ukur tanpa atom
yang sama $(X,\Sigma,\mu)$, maka
$\{x:x\in X$, $i\mapsto x(i)$ injektif$\}$ adalah
koterabaikan terhadap $\lambda$.

\spheader 251Wp Sekarang misalkan kita mempunyai keluarga lain
$\familyiI{(Y_i,\Tau_i,\nu_i)}$ ruang ukur, dengan produk
$(Y,\Lambda',\lambda')$, dan fungsi-fungsi \imp\ $f_i:X_i\to Y_i$ untuk setiap $i$;
definisikan $f:X\to Y$ dengan menyatakan bahwa $f(x)(i)=f_i(x(i))$ untuk $x\in X$ dan
$i\in I$. Jika semua $\nu_i$ $\sigma$-hingga, maka $f$ merupakan \imp\ bagi
$\lambda$ dan $\lambda'$.

\exercises{
\leader{251X}{Latihan dasar (a)}
%\spheader 251Xa
Misalkan $X$ dan $Y$ himpunan, $\Cal A\subseteq\Cal PX$ dan
$\Cal B\subseteq\Cal PY$. Misalkan $\Sigma$ $\sigma$-aljabar subhimpunan
$X$ yang dibangkitkan oleh $\Cal A$ dan $\Tau$ $\sigma$-aljabar subhimpunan
$Y$ yang dibangkitkan oleh $\Cal B$. Tunjukkan bahwa $\Sigma\tensorhat\Tau$ adalah
$\sigma$-aljabar subhimpunan $X\times Y$ yang dibangkitkan oleh
$\{A\times B:A\in\Cal A$, $B\in\Cal B\}$.
%251D query out of order

\spheader 251Xb Misalkan $(X,\Sigma,\mu)$ dan
$(Y,\Tau,\nu)$
ruang ukur; misalkan $\lambda_0$ ukuran produk primitif pada
$X\times Y$, dan $\lambda$ ukuran produk c.l.d.\ Tunjukkan bahwa
$\lambda_0W<\infty$ jika dan hanya jika $\lambda W<\infty$ dan $W$ termuat dalam suatu himpunan
berbentuk

\Centerline{$(E\times Y)\cup (X\times F)\cup\bigcup_{n\in\Bbb
N}E_n\times F_n$}

\noindent dengan $\mu E=\nu F=0$ dan $\mu E_n<\infty$, $\nu F_n<\infty$
untuk setiap $n$.
%251F

\sqheader 251Xc Tunjukkan bahwa jika $X$ dan $Y$ sebarang himpunan, masing-masing
dengan ukuran cacahnya, maka ukuran produk primitif dan c.l.d.\
pada $X\times Y$ keduanya merupakan ukuran cacah pada $X\times
Y$.
%251F

\spheader 251Xd Misalkan $(X,\Sigma,\mu)$ dan $(Y,\Tau,\nu)$
ruang ukur; misalkan $\lambda_0$ ukuran produk primitif pada
$X\times Y$, dan $\lambda$ ukuran produk c.l.d.\ Tunjukkan bahwa

$$\eqalign{\lambda_0&\text{ ditentukan secara lokal}\cr
&\iff\,\lambda_0\text{ semihingga}\cr
&\iff\,\lambda_0=\lambda\cr
&\iff\,\lambda_0\text{ dan }\lambda
  \text{ mempunyai himpunan terabaikan yang sama}.\cr}$$

\sqheader 251Xe (Lihat 251W.) Misalkan
$\langle(X_i,\Sigma_i,\mu_i)\rangle_{i\in
I}$ keluarga ruang ukur, dengan $I$ himpunan berhingga tak kosong.
Tetapkan $X=\prod_{i\in I}X_i$. Untuk $A\subseteq X$, tetapkan

\Centerline{$\theta A
=\inf\{\sum_{n=0}^{\infty}\prod_{i\in I}\mu_iE_{ni}:
E_{ni}\in\Sigma_i\Forall n\in\Bbb N,\,i\in
I,\,A\subseteq\bigcup_{n\in\Bbb N}\prod_{i\in I}E_{ni}\}$.}

\noindent Tunjukkan bahwa $\theta$ merupakan ukuran luar pada $X$. Misalkan
$\lambda_0$ ukuran yang didefinisikan dari $\theta$ dengan metode \Caratheodory,
dan untuk $W\in\dom\lambda_0$ tetapkan

\Centerline{$\lambda W=\sup\{\lambda_0(W\cap\prod_{i\in I}E_i):
E_i\in\Sigma_i,\,\mu_iE_i<\infty$ untuk setiap $i\in I\}$.}

\noindent Tunjukkan bahwa $\lambda$ merupakan ukuran pada $X$, dan merupakan versi c.l.d.\
dari $\lambda_0$.
%251I

\sqheader 251Xf (Lihat 251W.) Misalkan $I$ himpunan berhingga tak kosong dan
$\langle(X_i,\Sigma_i,\mu_i)\rangle_{i\in I}$ keluarga ruang
ukur. Untuk $K\subseteq I$ tak kosong, tetapkan $X^{(K)}=\prod_{i\in K}X_i$
dan misalkan $\lambda_0^{(K)}$, $\lambda^{(K)}$ ukuran-ukuran pada $X^{(K)}$
yang dikonstruksi seperti dalam 251Xe. Tunjukkan bahwa jika $K$ subhimpunan sejati
tak kosong dari $I$, maka bijeksi alamiah antara $X^{(I)}$ dan
$X^{(K)}\times X^{(I\setminus K)}$ mengidentifikasikan $\lambda_0^{(I)}$ dengan
ukuran produk primitif dari $\lambda_0^{(K)}$ dan
$\lambda_0^{(I\setminus K)}$, serta $\lambda^{(I)}$ dengan ukuran produk
c.l.d.\ dari $\lambda^{(K)}$ dan $\lambda^{(I\setminus K)}$.
%251Xe, 251I

\sqheader 251Xg Dengan menggunakan 251Xe-251Xf di atas, atau cara lain, tunjukkan bahwa jika
$(X_1,\Sigma_1,\mu_1)$, $(X_2,\Sigma_2,\mu_2)$,
$(X_3,\Sigma_3,\mu_3)$ ruang-ruang ukur, maka ukuran produk primitif dan c.l.d.\
$\lambda_0$, $\lambda$ pada $(X_1\times X_2)\times X_3$, yang dikonstruksi dengan
terlebih dahulu mengambil ukuran produk yang sesuai pada $X_1\times X_2$ lalu
mengambil produknya dengan ukuran pada $X_3$, diidentifikasikan
dengan ukuran-ukuran produk yang bersesuaian pada $X_1\times(X_2\times X_3)$ oleh
bijeksi kanonik antara himpunan $(X_1\times X_2)\times X_3$ dan
$X_1\times(X_2\times X_3)$.
%251Xf, 251Xe, 251I

\spheader 251Xh(i) Apa yang terjadi dalam 251Xe ketika $I$ merupakan
himpunan satu anggota? (ii) Rancang konvensi yang sesuai agar 251Xe-251Xf
tetap berlaku ketika satu atau lebih dari himpunan $I$, $K$, $I\setminus K$ di sana
kosong.
%251Xf, 251Xe, 251I

\sqheader 251Xi Misalkan $(X,\Sigma,\mu)$ ruang ukur lengkap yang
ditentukan secara lokal, dan $I$ sebarang himpunan tak kosong; misalkan $\nu$
ukuran cacah pada $I$. Tunjukkan bahwa ukuran produk c.l.d.\ pada
$X\times I$ sama dengan (atau setidaknya dapat diidentifikasikan dengan) ukuran
jumlah langsung dari keluarga $\langle(X_i,\Sigma_i,\mu_i)\rangle_{i\in I}$,
jika kita menetapkan $(X_i,\Sigma_i,\mu_i)=(X,\Sigma,\mu)$ untuk setiap $i$.
%251I

\sqheader 251Xj Misalkan
$\langle(X_i,\Sigma_i,\mu_i)\rangle_{i\in I}$ keluarga ruang
ukur, dengan jumlah langsung $(X,\Sigma,\mu)$
(214L). Misalkan $(Y,\Tau,\nu)$ sebarang ruang ukur, dan berikan $X\times Y$,
$X_i\times Y$ ukuran produk c.l.d.\ masing-masing. Tunjukkan bahwa bijeksi
alamiah antara $X\times Y$ dan
$Z=\bigcup_{i\in I}(X_i\times Y)\times\{i\}$
merupakan isomorfisme antara ukuran pada $X\times Y$ dan
ukuran jumlah langsung pada $Z$.
%251Xi, 251I

\sqheader 251Xk Misalkan $(X,\Sigma,\mu)$ sebarang ruang ukur, dan
$Y$ himpunan satu anggota
$\{y\}$; misalkan $\nu$ ukuran pada $Y$ sedemikian sehingga $\nu Y=1$.
Tunjukkan bahwa bijeksi alamiah antara $X\times\{y\}$ dan $X$
mengidentifikasikan ukuran produk primitif pada $X\times\{y\}$ dengan
$\check\mu$ sebagaimana didefinisikan dalam 213Xa, dan ukuran produk c.l.d.\ dengan
versi c.l.d.\
dari $\mu$. Jelaskan bagaimana menggabungkan hal ini dengan 251Xg dan
251Ic untuk membuktikan 251T.
%251I

\sqheader 251Xl Misalkan $(X,\Sigma,\mu)$ dan $(Y,\Tau,\nu)$ ruang
ukur, dan $\lambda$ ukuran produk c.l.d.\ pada $X\times Y$. Tunjukkan
bahwa $\lambda$ merupakan versi c.l.d.\ dari pembatasannya pada
$\Sigma\tensorhat\Tau$.
%251K

\spheader 251Xm Misalkan $(X,\Sigma,\mu)$ dan $(Y,\Tau,\nu)$ ruang
ukur, dengan ukuran produk primitif dan c.l.d.\ $\lambda_0$,
$\lambda$. Misalkan
$\lambda_1$ sebarang ukuran dengan domain $\Sigma\tensorhat\Tau$ sedemikian sehingga
$\lambda_1(E\times F)=\mu E\cdot\nu F$ kapan pun $E\in\Sigma$ dan
$F\in\Tau$. Tunjukkan bahwa $\lambda W\le\lambda_1W\le\lambda_0W$ untuk setiap
$W\in\Sigma\tensorhat\Tau$.
%251K

\spheader 251Xn Misalkan $(X,\Sigma,\mu)$ dan $(Y,\Tau,\nu)$ dua ruang
ukur, dan $\lambda_0$ ukuran produk primitif pada $X\times Y$.
Tunjukkan bahwa ukuran luar $\lambda_0^*$ yang bersesuaian tidak lain adalah
ukuran luar $\theta$ dari 251A.
%251P

\spheader 251Xo Misalkan $(X,\Sigma,\mu)$ dan $(Y,\Tau,\nu)$
ruang ukur, dan $A\subseteq X$, $B\subseteq
Y$ subhimpunan; tuliskan $\mu_A$, $\nu_B$ untuk ukuran-ukuran
subruang. Misalkan $\lambda_0$ ukuran produk primitif
pada $X\times Y$, dan $\lambda_0^{\#}$
ukuran subruang yang diinduksinya pada $A\times B$. Misalkan $\tilde\lambda_0$
ukuran produk primitif dari
$\mu_A$ dan $\nu_B$ pada $A\times B$. Tunjukkan bahwa $\tilde\lambda_0$
memperluas $\lambda_0^{\#}$.
Tunjukkan bahwa jika
{\it salah satu} ($\alpha$) $A\in\Sigma$ dan $B\in\Tau$
{\it atau} ($\beta$) $A$ dan $B$ keduanya dapat ditutupi oleh barisan himpunan
berukuran hingga
{\it atau} ($\gamma$) $\mu$ dan $\nu$ keduanya dapat dilokalkan secara ketat,
maka $\tilde\lambda_0=\lambda_0^{\#}$.
%251Q

\spheader 251Xp Dalam 251Q, tunjukkan bahwa $\tilde\lambda$ dan $\lambda^{\#}$
akan mempunyai ideal nol yang sama, sekalipun tidak satu pun syarat
251Q(ii) dipenuhi.
%251Q query out of order

\spheader 251Xq Misalkan $(X,\Sigma,\mu)$ dan $(Y,\Tau,\nu)$ sebarang ruang
ukur, dan $\lambda_0$ ukuran produk primitif pada $X\times Y$.
Tunjukkan bahwa $\lambda_0^*(A\times B)=\mu^*A\cdot\nu^*B$ untuk sebarang
$A\subseteq X$ dan $B\subseteq Y$.
%251S

\spheader 251Xr Misalkan $(X,\Sigma,\mu)$ dan $(Y,\Tau,\nu)$
ruang ukur, dan $\hat\mu$ pelengkapan dari $\mu$.
Tunjukkan bahwa pasangan $\mu$, $\nu$ dan $\hat\mu$, $\nu$ mempunyai ukuran produk
primitif yang sama.
%251T

\spheader 251Xs Misalkan $(X,\Sigma,\mu)$ ruang ukur semihingga.
Tunjukkan bahwa $\mu$ tanpa atom jika dan hanya jika diagonal $\{(x,x):x\in X\}$
terabaikan terhadap ukuran produk c.l.d.\ pada $X\times X$.
%251U

\spheader 251Xt Misalkan $(X,\Sigma,\mu)$ ruang ukur tanpa atom,
dan $(Y,\Tau,\nu)$ sebarang ruang ukur. Tunjukkan bahwa ukuran produk c.l.d.\
pada $X\times Y$ tanpa atom.
%251+

\sqheader 251Xu Misalkan $(X,\Sigma,\mu)$ dan $(Y,\Tau,\nu)$ ruang
ukur, dan $\lambda$ ukuran produk c.l.d.\ pada $X\times Y$.
(i) Tunjukkan bahwa jika $\mu$ dan $\nu$ atomik murni, maka $\lambda$ juga demikian.
(ii) Tunjukkan bahwa jika $\mu$ dan $\nu$ bertumpu pada titik, maka $\lambda$ juga demikian.
%251+  (ii) used in 491F

\leader{251Y}{Latihan lanjutan (a)}
%\spheader 251Ya
Misalkan $X$, $Y$ himpunan dengan $\sigma$-aljabar subhimpunan $\Sigma$, $\Tau$.
Andaikan $h:X\times Y\to\Bbb R$ terukur terhadap $\Sigma\tensorhat\Tau$
dan $\phi:X\to Y$ terukur terhadap $(\Sigma,\Tau)$ (121Yc). Tunjukkan bahwa
$x\mapsto h(x,\phi(x)):X\to\Bbb R$ terukur terhadap $\Sigma$.
%251D

\spheader 251Yb Tunjukkan bahwa terdapat
ruang-ruang ukur $(X_1,\Sigma_1,\mu_1)$ dan $(X_2,\Sigma_2,\mu_2)$,
ruang probabilitas $(Y,\Tau,\nu)$ dan fungsi \imp\ $f:X_1\to X_2$
sedemikian sehingga $h:X_1\times Y\to X_2\times Y$ bukan \imp\ bagi ukuran produk
c.l.d.\ pada $X_1\times Y$ dan $X_2\times Y$, dengan
$h(x,y)=(f(x),y)$ untuk $x\in X_1$ dan $y\in Y$.
%251L

\spheader 251Yc Misalkan $(X,\Sigma,\mu)$ ruang ukur lengkap yang
ditentukan secara lokal dengan subruang $A$ yang ukurannya tidak ditentukan secara lokal
(lihat 216Xb). Tetapkan $Y=\{0\}$, $\nu Y=1$ dan
pertimbangkan ukuran produk c.l.d.\ pada $X\times Y$ dan $A\times Y$;
tuliskan $\Lambda$, $\tilde\Lambda$ untuk domain masing-masing. Tunjukkan bahwa
$\tilde\Lambda$ memuat secara sejati $\{W\cap(A\times Y):W\in\Lambda\}$.
%251P

\spheader 251Yd Misalkan $(X,\Sigma,\mu)$ sebarang ruang ukur,
$(Y,\Tau,\nu)$ ruang ukur tanpa atom, dan $f:X\to Y$ fungsi
terukur terhadap $(\Sigma,\Tau)$. Tunjukkan bahwa $\{(x,f(x)):x\in X\}$
terabaikan terhadap ukuran produk c.l.d.\ pada $X\times Y$.
%251U
}%end of exercises

\endnotes{
\Notesheader{251} Terdapat kesulitan nyata dalam memutuskan konstruksi mana
yang harus dinyatakan sebagai produk `yang' baku dari dua ukuran sebarang.
Frasa saya `ukuran produk primitif', dan
notasi $\lambda_0$, mengungkapkan suatu kecenderungan; pilihan saya sendiri adalah
produk c.l.d.\ $\lambda$, karena dua alasan utama. Pertama,
$\lambda_0$ kemungkinan besar bersifat `buruk', khususnya,
tidak semihingga, bahkan jika $\mu$ dan $\nu$ bersifat `baik' (251Xd, 252Yk),
sedangkan
$\lambda$ mewarisi beberapa sifat terpenting dari $\mu$ dan
$\nu$ (misalnya, 251O); kedua, dalam hal ruang ukur topologis
$X$ dan $Y$, sering kali terdapat ukuran topologis kanonik
pada $X\times Y$, yang kemungkinan lebih erat berkaitan dengan
$\lambda$ daripada dengan $\lambda_0$. Namun, untuk menjelaskan hal ini saya
harus meminta Anda menunggu hingga \S417 dalam Jilid 4.

Ukuran produk `primitif' dapat saja dihapus sama sekali
dari uraian, atau setidaknya ditempatkan dalam latihan. Memang inilah
yang saya harapkan akan saya lakukan dalam bagian selanjutnya dari risalah ini, sebab (menurut
pandangan saya) semua ciri penting ukuran produk pada sejumlah berhingga
faktor dapat dinyatakan dalam istilah ukuran produk c.l.d. Namun,
untuk pengantar pertama mengenai ukuran produk, pendekatan langsung
ke ukuran produk c.l.d.\ (melalui deskripsi $\lambda^*$
dalam 251P, misalnya) terasa terlalu berat, dan saya memiliki semacam
kewajiban untuk menyajikan pesaing paling alamiah dari ukuran produk c.l.d.\
dengan cukup menonjol agar Anda dapat menilai sendiri apakah saya
benar ketika mengesampingkannya. Tentu ada hasil-hasil yang berkaitan dengan
ukuran produk primitif (251Xn, 251Xq, 252Yc) yang memiliki
kesederhanaan yang menarik.

Pertentangan itu tentu saja dapat dihindari sepenuhnya jika kita langsung mengkhususkan
diri pada ruang $\sigma$-hingga, yang pada ruang-ruang tersebut kedua
konstruksi bertepatan (251K). Namun, bahkan hal ini tidak menyelesaikan semua
masalah. Ada ukuran alternatif yang lazim, sering disebut ukuran produk `yang' baku:
pembatasan $\lambda_{0\Cal B}$ dari $\lambda_0$ pada
$\sigma$-aljabar $\Sigma\tensorhat\Tau$. (Lihat, misalnya,
{\smc Halmos 50}.) Keuntungannya ialah bahwa jika suatu fungsi $f$ pada
$X\times Y$
terukur terhadap $\Sigma\tensorhat\Tau$, maka $x\mapsto f(x,y)$ terukur terhadap
$\Sigma$ untuk setiap $y\in Y$. (Hal ini karena

$$\{W:W\subseteq X\times Y,\,\{x:(x,y)\in
W\}\in\Sigma\,\Forall \,y\in Y\}$$

\noindent merupakan $\sigma$-aljabar subhimpunan $X\times Y$ yang memuat
$E\times F$ kapan pun $E\in\Sigma$ dan $F\in\Tau$, sehingga memuat
$\Sigma\tensorhat\Tau$.) Keberatan utama, menurut saya, ialah bahwa
ukuran Lebesgue pada $\BbbR^2$ tidak lagi merupakan produk `yang' baku dari ukuran Lebesgue
pada $\Bbb R$ dengan dirinya sendiri. Secara umum, sudah selayaknya kita mencari
ukuran yang mengukur sebanyak mungkin himpunan, dan saya lebih suka menghadapi
masalah-masalah teknis (yang saya akui dapat menggentarkan) dengan
mencari definisi-definisi yang sesuai ketika mendekati teorema-teorema utama,
daripada mengandalkan perbaikan ad hoc ketika tiba saatnya menerapkannya.

Untuk sementara saya tidak memberikan contoh ukuran produk lebih lanjut,
karena penyelidikan contoh-contoh tertentu akan jauh
lebih mudah dengan bantuan hasil-hasil dari bagian berikutnya. Tentu saja contoh
utama, dan yang seharusnya selalu terlintas dalam pikiran ketika mendengar
kata-kata `ukuran produk', adalah ukuran Lebesgue pada
$\BbbR^2$, yaitu kasus $r=s=1$ dari 251N dan 251R. Untuk memperoleh gambaran
tentang apa yang dapat terjadi ketika
salah satu faktornya tidak $\sigma$-hingga, Anda dapat melihat lebih dahulu 252K.

Saya berharap Anda akan melihat bahwa definisi ukuran luar
$\theta$ dalam 251A bersesuaian dengan definisi baku ukuran luar Lebesgue,
dengan `persegi panjang terukur' $E\times F$ menggantikan
interval, dan fungsional $E\times F\mapsto\mu E\cdot\nu F$
menggantikan `panjang' atau `volume' suatu interval;
lebih jauh, jika $E$ dan $F$ dipandang sebagai interval, terdapat hubungan yang jelas
antara ukuran Lebesgue pada $\BbbR^2$ dan ukuran produk
pada $\Bbb R\times\Bbb R$. Tentu saja suatu `hubungan yang jelas' bukanlah
hal yang sama dengan teorema yang benar-benar dirumuskan dengan hipotesis dan
kesimpulan yang tepat, tetapi Teorema 251N jelas bersifat sentral. Namun, jauh sebelum itu,
terdapat kesejajaran lain antara konstruksi 251A dan
konstruksi ukuran Lebesgue. Dalam kedua kasus, bukti bahwa kita memperoleh
ukuran luar langsung mengikuti rumus definisinya (dalam 113Yd saya memberikan
sebagai latihan suatu hasil umum yang mencakup 251B), dan oleh karena itu suatu
konstruksi yang sangat umum dapat menuntun kita kepada suatu ukuran. Namun ukuran itu
akan jauh kurang menarik dan kurang bernilai kecuali jika ia membuat himpunan-himpunan
dasar, yang dalam hal ini adalah persegi panjang terukur, menjadi terukur dan memberikan
nilai ukuran yang benar kepada himpunan-himpunan tersebut.
Jadi 251E bersesuaian dengan teorema bahwa interval terukur terhadap Lebesgue,
dengan ukuran yang benar (114Db, 114G). Inilah kunci sesungguhnya
bagi konstruksi tersebut, dan merupakan salah satu gagasan mendasar dalam teori ukur.

Kesejajaran lain lagi terdapat dalam 251Xn; ukuran luar yang mendefinisikan
ukuran produk primitif $\lambda_0$ tepat sama dengan ukuran luar
yang didefinisikan dari $\lambda_0$. Saya menguraikan fenomena yang bersesuaian
untuk ukuran Lebesgue dalam 132C.

Setiap konstruksi yang mengklaim sebutan `kanonik' harus
memenuhi berbagai persyaratan alamiah; misalnya, kita mengharapkan
bijeksi kanonik antara $X\times Y$ dan $Y\times X$ merupakan suatu
isomorfisme antara ruang-ruang ukuran produk yang bersesuaian.
`Komutativitas' produk dalam pengertian ini, menurut saya, jelas dari
definisi-definisi dalam 251A-251C. Jelas pula bahwa kita menginginkan -- meskipun, menurut
saya, tidak jelas bahwa hal itu benar -- agar produk tersebut `asosiatif' dalam arti
bahwa bijeksi kanonik antara $(X\times Y)\times Z$ dan
$X\times(Y\times Z)$ juga merupakan suatu isomorfisme antara
produk-produk ukuran produk yang bersesuaian. Hal ini memang berlaku bagi
ukuran produk primitif maupun c.l.d.\ (251Wh, 251Xe-251Xg).

Dengan menelusuri klasifikasi ruang ukur yang disajikan dalam
\S211, kita mendapati bahwa ukuran produk primitif $\lambda_0$ lengkap
untuk sebarang ukuran faktor $\mu$, $\nu$, sedangkan ukuran produk c.l.d.\
$\lambda$ selalu lengkap dan ditentukan secara lokal.
$\lambda_0$ mungkin tidak
semihingga, bahkan jika $\mu$ dan $\nu$ dapat dilokalkan secara ketat
(252Yk); tetapi $\lambda$ akan dapat dilokalkan secara ketat jika $\mu$ dan
$\nu$ demikian (251O). Tentu saja hal ini berkaitan dengan kenyataan bahwa
ukuran produk c.l.d.\ bersifat distributif terhadap jumlah langsung (251Xj). Jika
$\mu$ atau $\nu$ tanpa atom, maka $\lambda$ juga demikian (251Xt). Baik
$\lambda$ maupun $\lambda_0$ bersifat $\sigma$-hingga jika $\mu$
dan $\nu$ demikian (251K). Mungkin saja $\mu$ dan $\nu$ keduanya
dapat dilokalkan tetapi $\lambda$ tidak (254U).

Setidaknya jika Anda telah mengerjakan Bab 21, kini Anda telah mempelajari cukup banyak
teori ukur `murni' sehingga penyelidikan semacam ini, betapapun
langsungnya, menimbulkan cukup banyak pertanyaan. Selain jumlah
langsung, kita juga mempunyai konstruksi `pelengkapan', `subruang', `ukuran
luar', dan (khususnya) `versi c.l.d.' untuk dipadukan dengan gagasan-gagasan
baru ini; saya memberikan beberapa hasil dalam 251T dan 251Xk. Mengenai
subruang, muncul beberapa kesulitan yang mungkin mengejutkan. Masalahnya ialah bahwa
ukuran produk pada produk dua subruang dapat mempunyai domain yang lebih luas
daripada yang mungkin diduga. Saya memberikan contoh sederhana dalam 251Yc dan contoh yang lebih
rumit dalam 254Yg. Untuk ruang yang dapat dilokalkan secara ketat, tidak ada
masalah (251Q); tetapi tampaknya tidak ada kriteria lain yang diambil dari daftar
sifat yang dipertimbangkan dalam \S251 yang memadai untuk meniadakan
kemungkinan munculnya fenomena yang meresahkan.

}%end of notes

\discrpage
